% mini-unison --- a formal model of Unison's content addressing and update scheme.
%
% THIS DOCUMENT IS FOR A READER WHO DOES NOT KNOW ARBOR, including Unison's own
% authors. Keep it that way:
%
%   * Sections 1-6 are about Unison alone. All comparison with arbor lives in
%     section 7, which introduces arbor for someone who has never heard of it.
%   * No repo-internal paths (docs/..., formalism/...), and no references to
%     arbor documents a reader cannot obtain.
%   * No process narration -- what an earlier draft got wrong, which prose read
%     missed what, which claims corrected which file. That material is real and
%     worth keeping, but it belongs in ../decisions.md (entry 2026-08-25), not
%     here. State findings positively; let the log carry the history.
%
% Source citations use short file keys defined in the table in section 1.7.
% Build: latexmk -pdf arbor-core.tex first (xr cross-refs), then this.

\documentclass[11pt]{article}

\usepackage[margin=1in]{geometry}
\usepackage{amsmath}
\usepackage{amssymb}
\usepackage{amsthm}
\usepackage{mathtools}
\usepackage{stmaryrd}
\usepackage{mathpartir}
\usepackage{xcolor}
\usepackage{xr-hyper}
\usepackage[colorlinks=true,linkcolor=blue!50!black,citecolor=green!40!black]{hyperref}
\usepackage[nameinlink]{cleveref}
\usepackage{enumitem}
\usepackage{array}
\usepackage{booktabs}

\externaldocument[core:]{arbor-core}
\externaldocument[stlc:]{arbor-stlc}

% Notation macros for arbor-core.
% Kept in one place so the eventual Agda names can be diffed against these.

% --- stores and configurations ---
\newcommand{\Store}{\Sigma}                 % term store
\newcommand{\Names}{N}                       % namespace
\newcommand{\Cache}{E}                        % evaluation cache
\newcommand{\Hist}{H}                         % binding history
\newcommand{\Config}[4]{\langle #1,\, #2,\, #3,\, #4\rangle}
\newcommand{\dom}{\operatorname{dom}}
\newcommand{\ran}{\operatorname{ran}}

% --- hashing ---
\newcommand{\Hash}{\mathbb{H}}               % the abstract type of hashes
\newcommand{\hashof}[1]{\lceil #1 \rceil}    % hash of a node
\newcommand{\Time}{\mathbb{T}}               % timestamps

% --- terms (deep, de Bruijn) ---
\newcommand{\Var}[1]{\mathsf{var}\,#1}
\newcommand{\Lam}[1]{\lambda.\,#1}
\newcommand{\App}[2]{#1\;#2}
\newcommand{\tRef}[1]{\mathsf{ref}\,#1}

% --- shallow nodes ---
\newcommand{\nVar}[1]{\mathsf{Var}(#1)}
\newcommand{\nLam}[1]{\mathsf{Lam}(#1)}
\newcommand{\nApp}[2]{\mathsf{App}(#1,#2)}
\newcommand{\nRef}[1]{\mathsf{Ref}(#1)}

% --- surface terms ---
\newcommand{\sVar}[1]{#1}
\newcommand{\sLam}[2]{\lambda #1.\,#2}
\newcommand{\sName}[1]{\mathsf{name}\,#1}

% --- de Bruijn operations (p4 ast.re) ---
\newcommand{\shift}[3]{\uparrow^{#1}_{#2}\!\left(#3\right)}   % shift by #1 above cutoff #2
\newcommand{\subst}[3]{[#1 \mapsto #2]\,#3}
\newcommand{\isclosed}[1]{\mathsf{closed}(#1)}

% --- store operations ---
\newcommand{\ingest}{\mathsf{ingest}}
\newcommand{\reconstruct}{\mathsf{reconstruct}}
\newcommand{\wf}{\mathsf{wf}}
\newcommand{\closedin}[2]{\mathsf{closed}_{#1}(#2)}  % #2 ∈ dom(#1) and reconstructs closed
\newcommand{\refs}{\mathsf{refs}}
\newcommand{\callers}{\mathsf{callers}}
\newcommand{\callersstar}{\mathsf{callers}^{*}}

% --- judgments ---
\newcommand{\evals}[3]{#1 \vdash #2 \Downarrow #3}          % Sigma |- t \Downarrow v
\newcommand{\resolves}[3]{#1 \vdash #2 \rightsquigarrow #3} % N |- x ~> h
\newcommand{\elab}[4]{#1 \vdash #2 \Rrightarrow #3}         % N |- s => t (elaboration)

% --- values ---
\newcommand{\val}{v}

% --- edit calculus ---
\newcommand{\steps}{\longrightarrow}
\newcommand{\bind}{\mathsf{bind}}
\newcommand{\rebind}{\mathsf{rebind}}
\newcommand{\unbind}{\mathsf{unbind}}
\newcommand{\multirebind}{\mathsf{multi\_rebind}}
\newcommand{\Pin}{\textsc{pin}}
\newcommand{\Follow}{\textsc{follow}}
\newcommand{\Explicit}{\textsc{explicit}}
\newcommand{\Migrate}{\mathsf{migrate}}
\newcommand{\substrefs}[2]{\{#1\}\,#2}       % ref-substitution application

% --- misc ---
\newcommand{\Name}{\mathsf{Name}}
\newcommand{\Node}{\mathsf{Node}}
\newcommand{\Term}{\mathsf{Term}}
\newcommand{\Val}{\mathsf{Val}}
\newcommand{\upd}[3]{#1[#2 \mapsto #3]}
\newcommand{\pfn}{\rightharpoonup}           % partial function arrow
\newcommand{\defeq}{\triangleq}

% Notation macros for mini-unison.
%
% DELIBERATELY DISJOINT from macros.tex. arbor's \Store, \Names, \Cache, \nRef
% and friends denote arbor's objects, whose TYPES differ here: a Unison store is
% Hash -> Component, not Hash -> Node, and a Unison namespace is a relation, not
% a partial function. Overloading the notation would hide exactly the differences
% this document exists to state, so every object gets a \U-prefixed name and
% renders with a subscript u.

% --- the three stored structures ---
\newcommand{\UStore}{\Sigma_{\mathsf{u}}}      % codebase: Hash -> Component
\newcommand{\UNames}{N_{\mathsf{u}}}           % namespace: Name -> P(Referent)
\newcommand{\UCache}{W}                        % watch cache: Hash -> Hash
\newcommand{\UConfig}[3]{\langle #1,\, #2,\, #3\rangle}

% --- references and components ---
\newcommand{\Pos}{\mathsf{Pos}}
\newcommand{\Comp}{\mathsf{Comp}}
\newcommand{\Referent}{\mathsf{Referent}}
\newcommand{\Reference}{\mathsf{Reference}}
\newcommand{\UDer}[2]{\mathsf{Derived}(#1,#2)}
\newcommand{\UBuiltin}[1]{\mathsf{Builtin}(#1)}
\newcommand{\URefT}[1]{\mathsf{Ref}(#1)}
\newcommand{\UCon}[2]{\mathsf{Con}(#1,#2)}

% --- object language (named ABT, NOT de Bruijn) ---
\newcommand{\uvar}[1]{#1}
\newcommand{\ulam}[2]{\lambda #1.\,#2}
\newcommand{\uapp}[2]{#1\;#2}
\newcommand{\unat}[1]{\underline{#1}}
\newcommand{\utext}[1]{\text{``}#1\text{''}}
\newcommand{\uref}[1]{\mathsf{r}\,#1}
\newcommand{\uann}[2]{#1 : #2}

% --- types ---
\newcommand{\TNat}{\mathsf{Nat}}
\newcommand{\TText}{\mathsf{Text}}
\newcommand{\uarr}[2]{#1 \mathbin{\to} #2}
\newcommand{\utypes}[3]{#1 \vdash #2 : #3}       % Gamma |- t : T

% --- hashing ---
\newcommand{\uhash}[2]{\lceil #2 \rceil_{#1}}    % hash in environment
\newcommand{\uhashc}[1]{\lceil #1 \rceil}        % closed hash
\newcommand{\HEnv}{\rho}                          % the de-Bruijn-at-hash-time environment
\newcommand{\phash}[1]{\mathsf{ph}(#1)}           % permutation hash
\newcommand{\nhash}[1]{\mathsf{nh}(#1)}           % name hash (second sort key)
\newcommand{\comphash}{\mathsf{hashComp}}
\newcommand{\hashcomps}{\mathsf{hashComps}}
\newcommand{\amb}{\mathsf{Ambiguous}}
\newcommand{\bg}{\mathcal{B}}                     % a binding group

% --- naming ---
\newcommand{\unconf}[2]{\mathsf{unconflicted}(#1,#2)}
\newcommand{\uresolve}[3]{#1 \vdash #2 \rightsquigarrow #3}

% --- update ---
\newcommand{\deps}{\mathsf{deps}}
\newcommand{\dependents}{\mathsf{dependents}}
\newcommand{\tdeps}{\mathsf{tdeps}}
\newcommand{\ppe}{\mathsf{ppe}}
\newcommand{\render}{\mathsf{render}}
\newcommand{\uparse}{\mathsf{parse}}
\newcommand{\uupdate}{\mathsf{update}}
\newcommand{\lib}{\mathsf{lib}}
\newcommand{\scratch}{\mathsf{scratch}}
\newcommand{\ok}{\mathsf{typechecks}}
\newcommand{\usteps}{\longrightarrow_{\mathsf{u}}}


\theoremstyle{definition}
\newtheorem{definition}{Definition}[section]
\newtheorem{example}{Example}[section]
\theoremstyle{plain}
\newtheorem{theorem}{Theorem}[section]
\newtheorem{lemma}{Lemma}[section]
\newtheorem{corollary}{Corollary}[section]
\newtheorem{proposition}{Proposition}[section]
\theoremstyle{remark}
\newtheorem{remark}{Remark}[section]

\newcommand{\src}[1]{{\footnotesize\textcolor{gray}{[\texttt{#1}]}}}
\newcommand{\arborcore}{\textsc{arbor-core}}
\newcommand{\arborstlc}{\textsc{arbor-stlc}}
\newcommand{\coreref}[1]{\arborcore~\cref{core:#1}}
\newcommand{\stlcref}[1]{\arborstlc~\cref{stlc:#1}}

\title{\textbf{mini-unison}\\[4pt]
\large A formal model of Unison's content addressing and update scheme}
\author{}
\date{Draft --- \today}

\begin{document}
\maketitle

\begin{abstract}
Unison is a programming language whose codebase is content-addressed. A definition is identified by
the hash of its syntax, names are metadata attached to hashes, and a computed result is cached
against a hash and never invalidated. These properties are widely cited. As far as we can determine
they have never been formally stated, and Unison has no peer-reviewed publication, so the
implementation is the only citable description of what it guarantees.

This document states them. Over a small object language, the simply-typed $\lambda$-calculus with
$\TNat$ and $\TText$, we model the following, faithfully to the implementation at a pinned commit:
hashing a strongly connected component of mutually recursive definitions; references as
$(\mathit{hash}, \mathit{position})$ pairs; reference transparency; the type inside the term hash;
the namespace as a relation in which conflicts are representable; the watch cache; \texttt{update},
which propagates an edit by rendering dependents to source, re-parsing and re-typechecking them; and
\texttt{merge}, which reconciles two independently edited namespaces by comparing a second hash
computed from names instead of content.

Four results are, we think, not obvious from the source. Two concern what a hash identifies: whether
the identity of a mutually recursive group can depend on what its members are \emph{named}, and
whether the interchangeability a hash confers applies to every member of such a group or only to one.
Two concern what happens when names are attached back onto hashes. \texttt{update} is correct only if
printing a definition and parsing it back preserves its hash, an obligation the implementation
maintains by testing and never states. \texttt{merge} distinguishes a human's edit from an automatic
rebuild by asking whether the definition \emph{looks} different, so a rename is indistinguishable
from authorship.

A closing section compares the design with \emph{arbor}, a content-addressed substrate the authors
are developing, for which a companion formalization exists. Arbor appears nowhere else.
\end{abstract}

\tableofcontents

% =====================================================================
\section{Introduction}\label{sec:uscope}

\subsection{What Unison does differently}

In a conventional language a definition is identified by its name and located in a file. In Unison a
definition is identified by the hash of its own syntax. Names, of its binders and of the definitions
it mentions, play no part in that identity. Three consequences follow:

\begin{itemize}[nosep]
  \item \textbf{Names are metadata.} Renaming a definition does not change it, so nothing is rebuilt
        and no cached result is invalidated. This holds of the machinery only: a rename is also the
        one action that can hand someone a merge conflict in code they never opened
        (\cref{sec:umerge}).
  \item \textbf{Definitions are immutable.} Editing one produces a new hash and so a new definition,
        and \texttt{update} rebuilds every dependent, transitively, into new definitions too
        (\cref{sec:uupdate}). The old ones still mean what they meant.
  \item \textbf{Caches need no invalidation.} A computed value is stored against the hash of the
        expression that produced it. That hash denotes one fixed expression forever, so the entry
        cannot go stale.
\end{itemize}

Mutually recursive definitions shape everything in \cref{sec:ustore}. A member of a cycle cannot
carry its neighbours' hashes, because computing any of them requires all of them. A cycle is
therefore hashed as a unit, with its internal references standing as indices. This is why the store
maps hashes to components, why a reference is a pair, and why \cref{def:comphash} needs a canonical
ordering.

Most of this document concerns the problems content addressing creates: names have to be reattached
to hashes when code is displayed, and propagating an edit to its dependents is a rebuild, not a
rename.

\subsection{Why formalize it}

Unison's documentation and blog posts state its claims about identity, immutability and caching as
design intent. No statement of them exists that is precise enough to check or to compare against
another design: what is identified with what, under which conditions, and what has to hold for the
conclusions to follow. Absent a publication, the implementation is the specification, and it is
roughly 161{,}000 lines of Haskell.

We wrote this to compare Unison with a system we are building (\cref{sec:contrast}). Several
observations below concern behaviour the code fixes but no prose description we have found states:
the position-$0$ boundary of \cref{prop:transparent}, the two orderings of \cref{rem:twoorders}, and
the three representations of \cref{rem:unhash}.

\subsection{Summary of results}\label{sec:summary}

Cross-references are for locating a statement later; this section reads without them.

\textbf{What a hash identifies.} Hashing a definition identifies it up to renaming of bound variables
and rotation of a recursive cycle, and no further. Two definitions that compute the same thing by
different means stay distinct. The identity of a recursive \emph{group} is only conditionally
independent of how it was written. The construction orders the group's members and uses their
positions as indices. When two members are indistinguishable except for which one they call, nothing
breaks the tie and the ordering falls back on what the members are \emph{named}. Unison detects this
at the top level and refuses to store such a hash. Inside a \texttt{let rec} the same detection runs
and its result is discarded, so the enclosing definition's identity depends on which binding was
typed first. The construction computes two different orderings of the same group: one fixes the
indices used while hashing, the other the positions stored in references. The interchangeability a
hash confers (substitute anything with the same hash, get the same enclosing hash) applies only to the
member at position $0$. Every other member is referenced opaquely, so reference transparency and the
group-hashing model are not independent features. A cyclic definition is not stored in the form that
was hashed: its self-references carry no hash, and an explicit inverse bridges the two forms.
\Cref{thm:ualpha,thm:ordering,cor:detect,rem:notpartial,rem:innercycle,def:comphash,rem:twoorders,prop:transparent,rem:unhash}.

\textbf{When a definition acquires its identity.} When the file containing it \emph{typechecks},
which happens on every save. Storing it and naming it are a separate, later gesture, so the
evaluation cache fills with no commit involved. A save also confers identities on things nothing ever
stores, so the cache can be keyed by a hash that denotes nothing in the codebase. The cache's
soundness argument never depended on the key being a definition.
\Cref{def:usave,rem:saveident,prop:watchdom}.

\textbf{Why the cache needs no invalidation.} A result is stored against the hash of the expression
that produced it, and that hash denotes one fixed expression forever. Staleness is unrepresentable.
This follows from immutability; it is not a policy that could be traded away. \Cref{cor:wnoinval}.

\textbf{What \texttt{update} rests on.} A stored definition refers to its dependencies by hash, so
making it call a new version means changing a hash inside it, and Unison never edits a stored tree.
It prints the dependent back out as source, where the dependency appears as a \emph{name}, and parses
the text again in a namespace where that name denotes the new version. The printer and parser perform
the substitution. \texttt{update} is therefore correct only if printing a definition and parsing it
back returns what it started from, apart from the reference meant to move. The implementation never
states this obligation; it maintains it as a golden-file regression corpus. Every named dependent is
rebuilt, and no caller can be left alone. \Cref{prop:print-parse,prop:follow-only}.

Substituting the new reference into each stored tree would look more direct but saves less than it
appears to. The dependents must be rebuilt, re-typechecked and rebound either way. The round trip
reuses the path ingest already provides, and its failure case comes free: when the rebuild does not
typecheck, the user gets back a file they can edit.

\textbf{What \texttt{merge} can and cannot tell.} Two codebases grown independently from a common
ancestor can be unioned without reconciling anything. This is the property content-addressed
codebases are usually sold on, and it is close to all that immutability buys. Every difficult case is
in the namespace, which is not content-addressed. There, for every definition a branch changed, merge
has to decide whether a person or a rebuild changed it. It prints the definition and asks whether it
reads as before. The test is a second hash over the printed form, with each reference replaced by the
name it is shown under; nothing stores it. A rebuilt definition still calls its dependency by the
same name, so it prints as before and passes. A renamed one prints differently, so a definition
nobody touched can be taken for someone's work.

Renaming is otherwise free: one row, no caller broken, no cached result lost, because no stored
definition contains a name. Merge is the single operation that reads names as content. A definition
nobody wrote cannot be conflicted provided neither party moved, aliased or deleted a name, and
changing a definition's names is the only user action that can hand someone a conflict in code they
never opened. It does not always arrive as a conflict. A rename carrying no edit produces none and
fails a step later, when the file is rendered under one branch's names and parsed against a namespace
holding the other's. A rename raced by an edit of the renamed definition merges cleanly and silently
undoes the rename. Unison computes the renames, but only those carrying no edit (the case that does
not conflict), and at the pinned commit passes the result to a debug printer. Merge is also not
commutative: the same two branches can merge cleanly in one direction and fail in the other, with
nothing reported either way.
\Cref{prop:merge-store,def:synhash,rem:syn-only,def:gestures,rem:renamefree,thm:noconflict,cor:noedit-user,rem:onlyrenames,prop:rename-conflict,cor:unchecked,prop:merge-asym,prop:merge-print-parse}.

\subsection{What this document is, and is not}

It is a model of an existing system, faithful to the implementation at commit \texttt{db60ce2}, and
it proposes no changes. Where we call something a cost or a limitation we mean it descriptively. The
one condition we identify that Unison treats as an error (\cref{cor:detect}) is already tracked
upstream as issue \#2787, and we think the implementation's response, detect and abort, is right.

We do not model Unison's language: there are no abilities, data declarations, pattern matching or
ANF runtime. \Cref{sec:unotmodeled} lists what is left out. Anything a Unison developer finds
surprising here is more likely our error than a discovery, and we welcome correction.

\subsection{The object language, and why it is typed}

The object language is the simply-typed $\lambda$-calculus with $\TNat$, $\TText$ and references.
It is typed for one reason. Unison's \texttt{update} re-typechecks the code it rebuilds, and when
that fails it removes the broken definitions from the namespace and hands them back as editable
source. Without a type system that branch is unreachable and cannot be stated.

\subsection{What is modeled}

Each operation a user can perform, what it is observed to do, and the gesture that performs it in
\texttt{ucm}:

\begin{itemize}[nosep]
  \item \textbf{Save.} You write definitions into a file. \texttt{ucm} watches the file and on each
        write typechecks all of it, computes an identity for every definition and watch expression,
        evaluates the watches and caches the results. Nothing is stored and no name is bound.
        \emph{In \texttt{ucm}: saving \texttt{scratch.u}, which is not a command.}
  \item \textbf{Commit, and propagate.} What the save identified is written to the codebase and the
        file's names are bound to it. Then everything that depended on a superseded version is
        rebuilt. Either all of it typechecks, or the broken definitions come back to you as source.
        The steps cannot be taken separately. \emph{In \texttt{ucm}: \texttt{update}, and
        \texttt{add}, which at this commit is the same command.}
  \item \textbf{Reference.} A stored definition may mention others. What it mentions is fixed when it
        is stored, so no edit anywhere alters what an existing definition means. \emph{In
        \texttt{ucm}: \texttt{view} shows a definition with names put back; \texttt{dependents}
        reads the relation the other way.}
  \item \textbf{Naming.} Names are attached to identities from outside and are not stored with them.
        A name may denote more than one thing at once, and the system continues to function when it
        does. \emph{In \texttt{ucm}: \texttt{ls} and \texttt{names} to look; \texttt{move.term},
        \texttt{alias.term} and \texttt{delete.term} to change. None touches a definition.}
  \item \textbf{Evaluation and caching.} A result is remembered against the identity of the
        expression that produced it and never has to be discarded. \emph{In \texttt{ucm}: a line of
        the file beginning \texttt{>}, evaluated by the save, its answer printed beneath it on every
        later save without being recomputed.}
  \item \textbf{Merge.} Two people edit from a common starting point and the results are combined.
        The system decides, definition by definition, whether each side's version was written by a
        person or produced by rebuilding, and asks only about the ones both people wrote. \emph{In
        \texttt{ucm}: \texttt{branch} to fork, \texttt{merge /other} to combine.}
\end{itemize}

Two things visible only through their consequences are also modeled, because the rest depends on
them: how the identity of a mutually recursive group is fixed (\cref{def:comphash}), and the fact
that a definition's recorded type is part of its identity (\cref{def:typeinhash}).

\subsection{Not modeled}\label{sec:unotmodeled}

\begin{itemize}[nosep]
  \item \textbf{The namespace's representation.} Unison's \texttt{Branch} is a hashed Merkle tree of
        name segments whose \texttt{children} field points at each child's \emph{causal} hash, so
        every level of the tree has its own history \src{branchh:22-39}. Here the namespace is a bare
        relation. \Cref{sec:unaming} models the naming semantics, which \texttt{update} needs, and
        not the object that carries them.
  \item \textbf{History and branching.} \texttt{Causal} as a content-addressed DAG, and with it the
        computation of two branches' common ancestor. \Cref{sec:umerge} models merge as a function
        of three namespaces and takes the ancestor as given.
  \item \textbf{Suffix resolution and the \texttt{lib}-depth priority rule.} Names here are atomic.
        In Unison a name is a sequence of segments and may be abbreviated to any unambiguous suffix.
  \item \textbf{Data declarations and \texttt{unique}-type minting.} With only $\TNat$, $\TText$ and
        arrows there are no declarations, so the $\UCon{\cdot}{\cdot}$ arm of $\Referent$ is present
        but vestigial and the minting axis is absent.
  \item \textbf{Abilities, the ANF runtime, \texttt{LocalIds}, and the serialization layer.}
\end{itemize}

% =====================================================================
\section{The object language}\label{sec:ulang}

There are two grammars. A term as \emph{written} mentions other definitions by name, and names are
variables. A term as \emph{stored} mentions them by reference, and a reference carries a hash. Name
resolution maps one to the other, and it is partial in a way that matters later.

In the implementation these are one datatype; the stage determines which constructors occur.
Splitting them here is an idealization that makes the rest of the document easier to read.

\subsection{Source terms}

\begin{definition}[Source terms]\label{def:uterm}
Fix a countable set of variables $x,y,z$.
\[
  e \;::=\; \uvar{x} \;\mid\; \ulam{x}{e} \;\mid\; \uapp{e}{e'}
        \;\mid\; \unat{n} \;\mid\; \utext{s} \;\mid\; \uann{e}{T}
\]
This is what a user writes. There is no way to write a hash, and no term former denotes a stored
definition. A mention of another definition is a \emph{free variable}, as a mention of a local
binding is a bound one, and only an environment tells the two apart.
\end{definition}

\subsection{References}

\begin{definition}[Reference]\label{def:uref}
\src{ref:77-119} A reference denotes something already stored.
\[
  r \;::=\; \UBuiltin{b} \;\mid\; \UDer{h}{i}
  \qquad b \in \TText,\; h \in \Hash,\; i \in \Pos = \mathbb{N}
\]
$\UBuiltin{b}$ names a primitive by text. $\UDer{h}{i}$ names the $i$-th member of the group stored
at $h$; \cref{def:ustore} explains why that is a pair.
\end{definition}

\begin{remark}[Primitives are named extrinsically]\label{rem:builtin}
An identifier is \emph{intrinsic} when it can be recomputed from the thing it names and
\emph{extrinsic} when it is an assertion by some authority. A hash is intrinsic; a builtin name is
extrinsic. The system is intrinsically identified everywhere except at its primitives, and the sum
type makes that boundary visible in the type of a reference. It is unavoidable: a hash chain must
bottom out at things whose content the system does not model.
\end{remark}

\subsection{Resolved terms}

\begin{definition}[Resolved terms]\label{def:urterm}
\[
  t \;::=\; \uvar{x} \;\mid\; \ulam{x}{t} \;\mid\; \uapp{t}{u}
        \;\mid\; \unat{n} \;\mid\; \utext{s} \;\mid\; \uref{r} \;\mid\; \uann{t}{T}
\]
\Cref{def:uterm} extended with references. Binders still carry names, and those names are stored;
\cref{def:uhash} erases them only while hashing.
\end{definition}

\begin{definition}[$\alpha$-equivalence]\label{def:ualpha}
$t \equiv_\alpha u$ is the least congruence identifying $\ulam{x}{t}$ with $\ulam{y}{t[x := y]}$ for
$y$ fresh in $t$. In the implementation it is the equality on terms, so $\alpha$-equivalent terms are
equal before any hashing takes place.
\end{definition}

\begin{remark}[Why storing binder names matters]\label{rem:names}
Storing de Bruijn indices and inventing names for display would produce the same hashes, so the
choice looks free. It is not. \texttt{update} (\cref{sec:uupdate}) renders stored definitions back to
source that a human reads and a parser resolves. Because names survive storage, the rendering uses
the author's variable names. A system that discarded them could state its round trip only up to a
choice of names.
\end{remark}

\begin{definition}[Name resolution]\label{def:uresolution}
\src{fparse:185-191; coreterm:172-176} Given a namespace $\UNames$ (\cref{def:unames}) and a set $B$ of
names bound by the enclosing binding group, resolution $\mathcal{R}^{B}_{\UNames}$ maps a source term
to a resolved term, homomorphically except at variables:
\[
  \mathcal{R}^{B}_{\UNames}(\uvar{x}) \;=\;
    \begin{cases}
      \uvar{x} & x \in B \text{, or } x \text{ locally bound}\\
      \uref{r} & \UNames(x) = \{\URefT{r}\}\\
      \uvar{x} & \text{otherwise}
    \end{cases}
\]
\end{definition}

\begin{remark}[Resolution is deliberately partial]\label{rem:staging}
Three of these cases leave a variable in place and only one produces a reference.

A name bound by the group being ingested cannot be resolved: its referent does not exist yet, and for
a recursive group it cannot exist until the group is hashed (\cref{sec:comphash}). Resolution stops at
the boundary of what is already stored, and ingest discharges the remaining variables.

A name that resolves to nothing, or to more than one thing, is left free on purpose. The
implementation's comments say so \src{coreterm:175-176}: ``not found: leave free for telling user
about expected type'' and ``ambiguous: leave free for TDNR''. An unresolved name survives into
typechecking, where the expected type may disambiguate it or produce a better error than resolution
could. Failure to resolve is not itself an error.
\end{remark}

\subsection{Types and typing}

\begin{definition}[Types and typing]\label{def:utyping}
\[
  T,U \;::=\; \TNat \;\mid\; \TText \;\mid\; \uarr{T}{U}
\]
Typing is over resolved terms. A context $\Gamma$ maps variables to types. A signature $\Delta$ maps
references to types: given for builtins, read off the store for derived references
(\cref{def:typeinhash}). Syntax-directed, with no inference:
\begin{mathpar}
  \inferrule*[right=U-Var]{\Gamma(x) = T}{\utypes{\Gamma}{\uvar{x}}{T}}
  \and
  \inferrule*[right=U-Nat]{\ }{\utypes{\Gamma}{\unat{n}}{\TNat}}
  \and
  \inferrule*[right=U-Text]{\ }{\utypes{\Gamma}{\utext{s}}{\TText}}
  \\
  \inferrule*[right=U-Abs]{\utypes{\Gamma, x{:}T}{t}{U}}{\utypes{\Gamma}{\ulam{x}{t}}{\uarr{T}{U}}}
  \and
  \inferrule*[right=U-App]{\utypes{\Gamma}{t}{\uarr{T}{U}} \\ \utypes{\Gamma}{u}{T}}
                          {\utypes{\Gamma}{\uapp{t}{u}}{U}}
  \\
  \inferrule*[right=U-Ref]{\Delta(r) = T}{\utypes{\Gamma}{\uref{r}}{T}}
  \and
  \inferrule*[right=U-Ann]{\utypes{\Gamma}{t}{T}}{\utypes{\Gamma}{\uann{t}{T}}{T}}
\end{mathpar}
A term with a free variable is not typeable. That is how \cref{rem:staging}'s unresolved names become
errors: at typechecking, with a type to report.
\end{definition}

\subsection{Binding groups}

\begin{definition}[Binding group]\label{def:bindgroup}
A \emph{binding group} is a finite map $\bg$ from names to resolved terms, whose members' free
variables all lie in $\dom(\bg)$. Ingest takes a binding group, not a single term.
\end{definition}

\begin{definition}[File]\label{def:ufile}
\src{ufile:270-290} A \emph{file} is a pair $F = (\bg, \overline{w})$ of a binding group and a
finite sequence of \emph{watch expressions} $\overline{w} = w_1,\dots,w_m$: resolved terms with no
name, whose free variables also lie in $\dom(\bg)$. The user edits a file, of which
\cref{def:bindgroup} describes only half.

A watch is anonymous, so nothing can refer to it and it is always its own singleton component. It is
closed over the group's names as a definition is, and \cref{def:uingest} hashes it in the same pass.
The implementation keeps the two together: \texttt{allTerms = join tlcs ++ join (snd <\$> watches)},
under a comment reading ``includes watches''.
\end{definition}

\begin{remark}[Why ingest is group-at-a-time]\label{rem:mutrec}
There is no recursive term former. Mutual recursion is several bindings that mention one another, so
the unit that can be given identities is the group. With single-term ingest every component would be
a singleton and \cref{def:comphash} would be invisible.

The closedness condition is what \cref{rem:staging} left over: the variables resolution could not
discharge are the group's own names, and ingest discharges them by hashing.
\end{remark}

% =====================================================================
\section{The store and the hash}\label{sec:ustore}

\subsection{The store and its components}

\begin{definition}[Component and store]\label{def:ustore}
A \emph{component} is a non-empty finite sequence of terms, $\Comp = \Term^{+}$: a strongly connected
component of the dependency graph, hashed as a unit. The store is
\[
  \UStore : \Hash \pfn \Comp .
\]
A definition is addressed by a \emph{pair}: $\UDer{h}{i}$ denotes the $i$-th member of the component
at $h$. A non-recursive definition is a singleton component with $i = 0$.
\end{definition}

\begin{remark}[The store is not hash-to-definition]\label{rem:storetype}
Content addressing is often summarized as ``a definition is identified by its hash'', and for
non-recursive code that is what happens. But the map is from a hash to a component, and a definition
is a pair. A member of a cycle cannot be hashed before its neighbours, since it refers to them and
they refer back. Hashing the cycle as a unit breaks the knot, at the price of making every reference
positional and every property below a property of components.
\end{remark}

\subsection{Hashing}

\begin{definition}[Token stream]\label{def:tokens}
\src{tokenizable:122-137} Nothing is hashed as raw bytes. The input is a tagged,
length-prefixed token stream: $\mathsf{Tag}$ is one byte; $\mathsf{Text}$ and $\mathsf{Bytes}$ are a
64-bit big-endian length followed by the payload; $\mathsf{Nat}$ is 64-bit big-endian;
$\mathsf{Hashed}$ is a fixed-width hash with no length prefix. Every stream begins with a
\emph{version} token, and each syntactic layer begins with a tag domain-separating it from the others
\src{term:147-149; type:136-139}. Write $\mathsf{acc}(\overline{\tau})$ for the hash of a token list.
The function is SHA3-512, untruncated.
\end{definition}

\begin{remark}[Versioning the hash function]\label{rem:version}
The leading version token makes a change of hash function change every hash, so simple values cannot
collide across generations in one global hash table. Underneath, the schema carries a
many-hashes-to-one-object indirection so that hashes computed under an old scheme keep resolving: ``if
the hashing algorithm for the object is changed'' \src{sql:33-49}. Few content-addressed systems plan
for this. The intended migration for a changed hash is versioning plus indirection, not
recomputation. It is provisioned but not exercised: every write hardcodes version $2$, with a
standing note to revisit ``once there are v1 and v2 hashes'' \src{queries:778}.
\end{remark}

\begin{definition}[Hashing a term]\label{def:uhash}
\src{abt:145-174} Hashing carries an environment
$\HEnv \in (\wp(\Name) + \Name)^{*}$, innermost first, where $\mathsf{R}(x)$ is an ordinary binder and
$\mathsf{L}(V)$ is a \emph{cycle frame} naming a whole set at once. Write $\uhash{\HEnv}{t}$.
\begin{align*}
  \uhash{\HEnv}{\uvar{x}}    &= \mathsf{acc}[\mathsf{Nat}\; j]
     && \text{$j$ least with $x \in \HEnv_j$; undefined if none}\\
  \uhash{\HEnv}{\ulam{x}{t}} &= \mathsf{acc}[\ldots,\ \mathsf{Hashed}\;\uhash{\mathsf{R}(x)\cdot\HEnv}{t}]\\
  \uhash{\HEnv}{\uapp{t}{u}} &= \mathsf{acc}[\ldots,\ \mathsf{Hashed}\;\uhash{\HEnv}{t},\ \mathsf{Hashed}\;\uhash{\HEnv}{u}]
\end{align*}
and correspondingly at the other formers. Three details matter. Binder names are discarded, so
$\alpha$-equivalent terms hash equally. Source annotations are excluded by construction: ``we ignore
annotations in the \texttt{Term}, as these should never affect the meaning of the term''
\src{abt:136-137}. Here ``annotation'' is the ABT's annotation parameter, which carries source spans;
a type ascription $\uann{t}{T}$ is a term former and is hashed, which is what makes
\cref{def:typeinhash} work. Finally, a variable bound by a cycle frame $\mathsf{L}(V)$ hashes to the
frame's index, the same index for every member of $V$, which \cref{def:comphash} exploits. A free
variable is an error.
\end{definition}

\begin{definition}[The type is inside the term hash]\label{def:typeinhash}
\src{term:92-109} Before hashing, each member of a binding group is wrapped in its type,
$t \mapsto \uann{t}{T}$, replacing an existing top-level annotation if there is one. The hash of a
definition is therefore a function of its type as well as its body, and a change of signature is a
change of identity.
\end{definition}

\begin{remark}[The exception that proves the rule]\label{rem:typeinhash}
A \texttt{hashTermComponentsWithoutTypes} \src{term:120-122} exists for use where no type is
available, and its caller is warned off: ``this shouldn't be used when storing terms in the codebase,
as it doesn't incorporate the type into the hash'' \src{convert:87-89}. Hashing without the type is
possible and is treated as the special case.
\end{remark}

\begin{remark}[\textsc{U-Ref} is a lookup, not a re-derivation]\label{rem:uref-delta}
Typing a reference consults $\Delta$ and requires no derivation about the target. This is affordable
because the type is inside the hash (\cref{def:typeinhash}). A stored definition's recorded type
cannot disagree with the definition, since changing the type would change the hash and produce a
different definition. The store admits no ill-typed entry, so \textsc{U-Ref} has nothing to check.
\end{remark}

\subsection{Reference transparency, and its boundary}

A hash is meant to be interchangeable with the thing it names: if two terms hash alike, putting one
where the other was should leave every enclosing hash unchanged. This holds here, for less of the
store than it appears to.

\begin{proposition}[Transparency holds exactly at position $0$]\label{prop:transparent}
\src{term:135-146} In the term layer,
\[
  \uhash{\HEnv}{\uref{\UDer{h}{0}}} = h,
  \qquad
  \uhash{\HEnv}{\uref{\UDer{h}{i}}} = \mathsf{acc}[\mathsf{Tag}\;1,\ \mathsf{Hashed}\;h,\ \mathsf{Nat}\;i]
  \quad (i \neq 0).
\]
\end{proposition}

\begin{remark}[What transparency does, and where its boundary falls]\label{rem:transparent}
For $i = 0$, which covers every non-recursive definition and one member of each cycle
(\cref{sec:comphash} assigns the positions), a reference contributes literally the referent's hash.
A reference to position $0$ is therefore interchangeable, inside any enclosing hash, with any term
that hashes to $h$.

For $i \neq 0$ it is opaque, tagged with its position. So transparency and the component model are
not independent features; the boundary between them falls at $\Pos = 0$. We do not know whether this
is intended (inlining a singleton component is free, which is a plausible reason) or an artifact of
the encoding. The source comment justifies only the $0$ case.

\Cref{ex:transparency} works through what this gives, which is less than the usual summary claims.
\end{remark}

\subsection{Component hashing}\label{sec:comphash}

This mechanism's properties are the least apparent from the code.

\begin{definition}[Component hashing]\label{def:comphash}
\src{abt:73-100, 175-217} Let $C = \langle (y_1,u_1),\dots,(y_k,u_k)\rangle$ be a
component, \emph{listed in name order}; the implementation always obtains it as
$\mathsf{Map.toList}$.\footnote{Strictly the order is on $(\mathit{freshening\ id}, \mathit{var\
kind})$, not the name. It coincides with alphabetical order here because top-level bindings have their
freshening id reset to $0$ and all carry the same var kind. The distinction matters on the re-hash
path of \cref{rem:unhash}, where the variables are machine-generated from stored positions and the
listing is therefore by $\Pos$. Re-hashing agrees with hashing because \cref{thm:ordering}(a)'s side
condition makes the two listings induce the same $\pi$.} Let $V = \{y_1,\dots,y_k\}$.

\emph{Step 1: permutation hashes.} In the environment $\HEnv_p = \mathsf{L}(V)$, every intra-cycle
reference resolves to the same index, so member bodies become comparable before any of them has an
identity:
\[
  \phash{y_j} \;=\; \uhash{\HEnv_p}{u_j}.
\]

\emph{Step 2: the binding order.} Sort the members by $\phash{\cdot}$, \textbf{stably}, giving a
permutation $\pi$; set $\HEnv = \mathsf{R}(y_{\pi(1)})\cdots\mathsf{R}(y_{\pi(k)})$. Each member now
has a distinct de Bruijn index, and the member hashes are $h_j = \uhash{\HEnv}{u_{\pi(j)}}$.

\emph{Step 3: the storage order.} With
$\overline{\tau} = \mathsf{Tag}\;1,\ \mathsf{Hashed}\;h_1,\dots,\mathsf{Hashed}\;h_k$ shared by every
member, put
\[
  \nhash{y} \;=\; \mathsf{acc}\bigl(\overline{\tau} \;+\!\!+\; [\mathsf{Hashed}\;\uhash{\HEnv}{\uvar{y}}]\bigr),
\]
so a member is disambiguated only by its own index in $\HEnv$. Sort the members by $\nhash{\cdot}$.
That order is the component's storage order, the position $i$ in $\UDer{h}{i}$ counts in it, and
\[
  \comphash(C) \;=\; \mathsf{acc}[\mathsf{Hashed}\;\nhash{y_{\sigma(1)}},\dots,\mathsf{Hashed}\;\nhash{y_{\sigma(k)}}]
\]
for $\sigma$ the sorting permutation.
\end{definition}

\begin{remark}[Two orders, not one]\label{rem:twoorders}
$\pi$ and $\sigma$ are different permutations of the same component. $\pi$ fixes the de Bruijn
indices used while hashing member bodies; $\sigma$ fixes the $\Pos$ indices that appear in stored
references. They coincide only by accident, since $\sigma$ sorts by a hash of the $\pi$-index.
Nothing observable depends on the difference, because re-hashing re-derives $\HEnv$ from scratch. A
reimplementation that conflated them would produce different hashes for the same code, and no prose
description of the algorithm we know of mentions that there are two.
\end{remark}

\begin{definition}[Ingest]\label{def:uingest}
\src{abt:106-134; ufile:274-284} $\hashcomps(F)$ decomposes $\bg \cup \overline{w}$ into strongly
connected components in topological order and folds \cref{def:comphash} over them, replacing the
names of already-hashed components with $\UDer{h}{i}$ before hashing the next. It returns a reference
for every member, definitions and watches alike.

$\hashcomps$ only computes references. It extends neither $\UStore$ nor $\UNames$; the two
transitions below do that, and the separation matters (\cref{def:usave}).
\end{definition}

\begin{definition}[Save]\label{def:usave}
\src{load:464-494; ufile:270-290} \emph{Saving} a file is the transition
\[
  \UConfig{\UStore}{\UNames}{\UCache}
  \;\longrightarrow\;
  \UConfig{\UStore}{\UNames}{\UCache'}
\]
given by: typecheck $F$; if it fails, stop. Otherwise compute $\hashcomps(F)$, and for each watch
$w_j$ with reference $r_j$, set $\UCache' = \UCache[r_j \mapsto v_j]$ where $v_j$ is $w_j$'s value,
evaluated against $\UStore$ together with the definitions $\hashcomps$ just gave references to, and
\textbf{only when $r_j \notin \dom(\UCache)$ already}.

$\UStore$ is unchanged: a save stores nothing. $\UNames$ is unchanged: a save names nothing. A save
confers identity and populates the cache.
\end{definition}

\begin{remark}[Identity is conferred by typechecking, not by committing]\label{rem:saveident}
A definition acquires its identity when the file typechecks, on every save, not when it is added. A
\texttt{TypecheckedUnisonFile} carries a reference $r$ (\cref{def:uref}) for every member in its
\texttt{hashTermsId} field, and nothing has been added to anything. \Cref{def:update} then stores and
binds what already has a hash. The commit gate of \cref{rem:mutrec} is easily misread as putting
identity at the \texttt{add}; it does not.

Two consequences follow. First, the gate is on identity, not storage, so a buffer that does not
typecheck has no identities at all, not merely nothing stored. Second, caching (\cref{sec:ucache})
happens at save time, so re-saving an unchanged expensive watch is free with no commit anywhere. This
is observable: on \texttt{ucm} release/1.1.1 one expensive watch costs 5.79\,s, the same watch saved
twice costs 5.82\,s, and two distinct watches of equal cost 11.45\,s.

The implementation makes hit and miss explicit: \texttt{evaluateWatches} returns an $\mathit{isHit}$
flag per watch, and the save path writes the cache ``only when there are no errors'' and
$\mathsf{not}\ \mathit{isHit}$ \src{load:490-493}.
\end{remark}

\begin{remark}[Where a watch parts company from a definition]\label{rem:watchsplit}
Both are hashed by the same call and treated alike up to \cref{def:usave}. They diverge at
\cref{def:update}: a definition's component is added to $\UStore$ and its name bound in $\UNames$; a
watch's is added to neither, and only its value survives, in $\UCache$ at the reference the save
computed. The system mints identities for things it never stores (\cref{prop:watchdom}).
\end{remark}

\begin{proposition}[The cache's domain is not contained in the store's]\label{prop:watchdom}
There are configurations reachable by \cref{def:update} in which
$\dom(\UCache) \not\subseteq \{\,\UDer{h}{i} \mid h \in \dom(\UStore)\,\}$.
\end{proposition}
\begin{proof}
Save (\cref{def:usave}) a file whose binding group is empty and whose watch list is a single closed
expression $w$. Then $w$ receives a reference $\UDer{h}{0}$ and $\UCache$ is extended at it, while
$\UStore$ is unchanged by definition. No commit is required, and none would help: \cref{def:update}
adds components for the file's definitions, of which there are none. Since $w$ is unnamed, no later
operation can reach it: $\tdeps$ ranges over named references (\cref{def:uscope}), and $\render$ has
no name to print it by (\cref{def:ppe}). So $h \notin \dom(\UStore)$ permanently.
\end{proof}

\begin{remark}[A cyclic definition is not stored in the form that was hashed]\label{rem:unhash}
An intra-component reference has three representations.

\emph{While hashing}, a reference from one cycle member to another is a \textbf{variable}, resolved
by \cref{def:uhash} to an index into the cycle frame. It must be, because the component's hash is
what is being computed.

\emph{In memory}, once \cref{def:uingest} has finished, those variables have been substituted by
$\UDer{h}{i}$ carrying the component's own freshly computed hash.

\emph{Serialized}, a self-reference is stored with \textbf{no hash at all}. The hash field becomes
optional, and the absent case marks a reference into the enclosing component. The type is named for
it, ``a possibly-self (R = \emph{recursive}) reference'' \src{ref:51-52}, and the storage format says
what absence means: ``an added distinguished element (here \texttt{Nothing}) which represents a
self-reference'' \src{tformat:19-23}.

So the bytes on disk are not the bytes that were hashed, and re-deriving a stored component's hash
requires undoing the substitution first. That operation is \texttt{unhashComponent}
\src{cterm:265-276}. Its declaration analogue is documented as ``replaces all \texttt{Nothing}
self-referential hashes with a variable reference to the relevant piece of the component''
\src{decl:81-83}, and the hash-validation path calls it before re-hashing and compares
\src{termhash:26-45}.

None of this is observable through the API, and the encoding is sensible: writing a component's own
hash into its members' serialized bytes would mean the bytes could not be produced until the hash was
known, which is the circularity the construction exists to break. We state it because ``the hash of
a definition is the hash of what is stored'' is a natural assumption, and for cyclic definitions it
fails in three ways.
\end{remark}

\subsection{What component hashing guarantees}

For one definition, identity is independent of presentation unconditionally, and \cref{thm:ualpha}
says how far that goes. For a recursive group it is conditional. \Cref{def:comphash} orders the
members by comparing their bodies. When two bodies are indistinguishable (for members of a cycle,
indistinguishable except for which neighbour each calls), the comparison cannot break the tie, and it
falls to what the members are called.

\begin{theorem}[Conditional presentation-independence]\label{thm:ordering}
Let $C$ and $C'$ be the same component listed in two different orders.
\begin{enumerate}[nosep,label=(\alph*)]
  \item If $\phash{\cdot}$ is \emph{injective} on $V$, then $\comphash(C) = \comphash(C')$, the member
        hashes agree, and each member receives the same $\Pos$.
  \item If $\phash{\cdot}$ is not injective, the tied members' de Bruijn indices are decided by the
        listed order, and hence, since the list is always obtained in name order, by the
        members' \textbf{names}. Two consequences, of different strengths:
        \begin{enumerate}[nosep,label=(b\arabic*)]
          \item \emph{Always}: which member receives which $\Pos$ depends on the names. Renaming a
                member can change what $\UDer{h}{0}$ denotes, and therefore which member of the cycle
                is referenced transparently (\cref{prop:transparent}).
          \item \emph{In general, but not for every tie}: $\comphash(C)$ itself depends on the names.
        \end{enumerate}
\end{enumerate}
\end{theorem}
\begin{proof}[Proof sketch]
$\HEnv_p$ mentions no member individually, so $\phash{y}$ is a function of $u_y$ alone and is
unaffected by the listing. (a) If $\phash{\cdot}$ is injective, sorting by it is a total order on $V$
determined by the members' bodies, so $\pi$, hence $\HEnv$, hence every $h_j$, hence every
$\nhash{\cdot}$, hence $\sigma$ and $\comphash$, are all listing-independent. (b) If
$\phash{y} = \phash{y'}$ for $y \neq y'$, the sort in Step 2 is stable and returns them in listed
order; exchanging them in the listing exchanges their de Bruijn indices. That gives (b1) immediately,
since $\HEnv$ is what $\nhash{\cdot}$, hence $\sigma$, hence $\Pos$, is computed from. It does not
give (b2) immediately, because the last step of \cref{def:comphash} accumulates the $\nhash{\cdot}$
values sorted: if exchanging two members merely permutes the multiset $\{\nhash{y}\}$, $\comphash$ is
unchanged and only the assignment moves. (b2) needs the exchange to change the multiset, which it
does whenever it changes some $h_j$. Since $h_j$ is computed in $\HEnv$, this is the typical case.
\Cref{ex:ambiguous} exhibits both behaviours.
\end{proof}

\begin{corollary}[The side condition is decidable, and is computed]\label{cor:detect}
\src{abt:186-192, 210-217} Injectivity of $\phash{\cdot}$ on $V$ is decided by grouping the
members by permutation hash. \texttt{doHashCycle} does this and emits
\texttt{IncompleteElementOrderingError}, listing the structurally equivalent sets.
\end{corollary}

\begin{remark}[Total function, conditional guarantee, aborting caller]\label{rem:notpartial}
$\comphash$ is a \emph{total} function: on the ambiguous input it returns a hash. What fails is
presentation-independence, the property content addressing exists to provide, and the fallback is
alphabetical order of user-chosen names. The usual summary (``canonical ordering is a partial
function'', ``Unison refuses'') is therefore wrong on both counts.

The refusal is at the caller. \texttt{doHashCycle} returns the condition as a warning, and its own
comment says ``hashing always succeeds even if it generates warnings''. \texttt{crashOnHashingWarning}
turns it into an exception, and the main term-component path takes that route \src{ufile:284}. A
top-level handler renders it as bug \texttt{E253299} with a workaround: add a distinguishing binding
such as \texttt{\_ = "this is the foo definition"} \src{abt:41, 55-62; main:116}.

For a top-level component, then, a decidable side condition guards presentation-independence, and the
implementation aborts instead of storing a hash that violates it. No top-level component ever receives
a presentation-dependent identity, which is stronger than a partial function. For a formalism the
obligation is sharper: the side condition has to appear in the statement of every downstream result
about component identity. The condition is tracked upstream as unisonweb/unison\#2787.

\Cref{rem:innercycle} describes the one caller for which none of this holds.
\end{remark}

\begin{remark}[Inside a \texttt{let rec} the same condition is computed and discarded]\label{rem:innercycle}
\src{abt:169-174, 202-206} A recursive group nested inside a definition is hashed by the same
\texttt{doHashCycle}, reached through a different caller:
\begin{quote}\small
``We ignore incomplete element ordering warnings when calling in from \texttt{hash'}; we don't want to
error on that when hashing internal let-bindings.''
\end{quote}
The warning is bound to \texttt{\_warnings} and dropped. Step 2 of \cref{def:comphash} sorts stably,
so tied members keep their listing order, which fixes their de Bruijn indices and so the enclosing
definition's hash. For an inner cycle the listing is not reset to name order as top-level bindings are
(\cref{def:comphash}, footnote), so $\pi$ inherits the order the bindings were written in.

Three definitions of the same function show it:
\[
\begin{array}{l@{\qquad}l@{\qquad}l}
  \texttt{m1 n = (a x = b x;\ b x = a x;\ a n)} & \texttt{\#2d5bs110vr} & \\
  \texttt{m2 n = (p x = q x;\ q x = p x;\ p n)} & \texttt{\#2d5bs110vr} & \just{members renamed} \\
  \texttt{m3 n = (b x = a x;\ a x = b x;\ a n)} & \texttt{\#qqn1r2c1od} & \just{members reordered}
\end{array}
\]
Renaming the two members changes nothing, as \cref{thm:ualpha} promises. Writing them in the other
order changes the identity of the containing definition, though no reader could tell the two apart.
Lifting the same pair of bindings to the top level aborts with \texttt{E253299}.

The guarantee of \cref{rem:notpartial} therefore concerns top-level components, not hashing in
general: one nesting level down, the system silently produces a presentation-dependent identity. The
dependence there is on source order, a different coordinate from \cref{thm:ordering}'s names and
arguably a worse one, since a name is at least something the user chose. The published workaround
(add a distinguishing binding) applies unchanged, but nothing reports that it is needed.
\end{remark}

\begin{theorem}[Hashing collapses $\alpha$-equivalence and cycle permutation]\label{thm:ualpha}
Let $\phash{\cdot}$ be injective on $V$. Then $\comphash$ identifies two components exactly when they
are related by $\alpha$-equivalence of members (\cref{def:ualpha}) together with a permutation of the
cycle. It does not identify observationally equivalent members.
\end{theorem}
\begin{proof}[Proof sketch]
$(\Leftarrow)$ Binder names are discarded by \cref{def:uhash}, and permutation-independence is
\cref{thm:ordering}(a). $(\Rightarrow)$ By injectivity of $\mathsf{acc}$ on token streams and the
domain-separating tags of \cref{def:tokens}, downward through the structure.
\end{proof}

% =====================================================================
\section{Naming}\label{sec:unaming}

Names are not part of a definition, and nothing above consulted one. This section says what names are
attached to and what a namespace is, which is what \cref{sec:uupdate} needs. The object that carries a
namespace, and its history, are out of scope (\cref{sec:unotmodeled}).

\begin{definition}[Referent]\label{def:ureferent}
\src{referent:27-30} $\Referent ::= \URefT{r} \mid \UCon{r}{k}$. Names denote
\emph{referents}, not references, because a name may denote a data constructor. With no data
declarations in scope (\cref{sec:unotmodeled}) the second arm is vestigial here, retained so that
\cref{def:unames} has the right shape.
\end{definition}

\begin{definition}[Namespace]\label{def:unames}
\src{branchh:22-39; sql:271} A namespace is a \emph{relation}
\[
  \UNames \;\subseteq\; \Name \times \Referent ,
\]
written $\UNames(x)$ for $\{\,\varrho \mid (x,\varrho) \in \UNames\,\}$. One name may denote several
referents and one referent may have several names; both are ordinary states, not errors. The derived
view has the same shape: names are recovered in memory as a relation between referents and names
\src{branchtype:110-111}. The database agreed, in the name-lookup index whose key included the
referent columns \src{sql:271}. That index was deleted in 2025 and the fragment is now vestigial (see
\cref{rem:conflict}).
\end{definition}

\begin{definition}[Unconflictedness]\label{def:unconf}
\src{branchtype:280-289} $\unconf{\UNames}{x}$ holds iff
$|\UNames(x)| = 1$, and $\unconf{\UNames}{}$ holds iff it holds at every name in $\dom(\UNames)$.
When it does, $\UNames$ determines a partial function and we write
$\uresolve{\UNames}{x}{\varrho}$. It is a derived view, recomputed and not maintained, and it is
all-or-nothing: \texttt{asUnconflicted} returns either the whole unconflicted namespace or the
offending name, never the unconflicted part of a namespace. So operations have $\unconf{\UNames}{}$
available as a predicate, not $\unconf{\UNames}{x}$.
\end{definition}

\begin{remark}[Conflict is representable; operations demand its absence]\label{rem:conflict}
Making conflict a legal state lets a merge complete without resolving everything first: the namespace
holds both candidates and the user sorts it out afterwards. The cost lands elsewhere. \texttt{update}
cannot proceed against a conflicted namespace and asserts unconflictedness before doing anything
\src{update:99-102}. Unison's forward-looking \texttt{BranchV3} format removes conflicted names
outright: ``a V3 branch is a trimmed-down V2 branch: Names can't be conflicted''
\src{branchv3:16-30}. The generality is real, and the direction of travel appears to be away from it.
\end{remark}

\begin{definition}[Namespace gestures]\label{def:gestures}
Besides \texttt{update} (\cref{def:update}), a user has three operations that change $\UNames$ and
nothing else. Writing $\varrho$ for the referent $\UNames$ binds at $x$:
\[
\begin{array}{l@{\qquad}l}
  \texttt{move}(x,y)   & \UNames' = \bigl(\UNames \setminus \{(x,\varrho)\}\bigr) \cup \{(y,\varrho)\} \\
  \texttt{alias}(x,y)  & \UNames' = \UNames \cup \{(y,\varrho)\} \\
  \texttt{delete}(x)   & \UNames' = \UNames \setminus \bigl(\{x\} \times \Referent\bigr)
\end{array}
\]
None of the three reads or writes $\UStore$, none can fail on a typing ground, and none changes any
hash. They are \texttt{move.term}, \texttt{alias.term} and \texttt{delete.term} in the interface.

The notation carries two preconditions. Writing $\varrho$ for ``the referent at $x$'' presupposes
$\unconf{\UNames}{x}$. And the target of a \texttt{move} or an \texttt{alias} must be unbound. The
relation could hold the union ($y$ would become conflicted), but the commands refuse with \emph{``A
term named \texttt{bar} already exists.''} Conflict remains representable and no gesture here can
produce one. This is a third instance of \cref{rem:conflict}'s pattern, after \cref{def:unconf}'s
all-or-nothing view and \cref{def:update}'s precondition.
\end{definition}

\begin{remark}[Renaming is free, and this is the usual reason given for content addressing]\label{rem:renamefree}
\Cref{def:gestures} is the formal content of the claim that a content-addressed codebase makes
renaming free. A rename is one row. It cannot break a caller or invalidate a cached result
(\cref{cor:wnoinval}), and no definition is rewritten, because no definition contains the name.
Nothing in \crefrange{sec:ustore}{sec:uupdate} can observe that one happened.

This stops being true in \cref{sec:umerge}. Merge is the one operation that reads names as content,
and \cref{cor:noedit-user} states the consequence.
\end{remark}

% =====================================================================
\section{The evaluation cache}\label{sec:ucache}

\begin{definition}[Watch cache]\label{def:watch}
\src{ctype:77-78; sql:121-127} The cache is a finite partial function
\[
  \UCache \;:\; (\Hash \times \Pos) \pfn \Term ,
\]
from the identity of an expression to its value. Two details are easy to get wrong. The key is a
\emph{reference}, not a bare hash: the primary key is $(\mathit{hash}, \mathit{component\ index})$, so
a watch is addressed the way \cref{def:uref} addresses a definition. The value is the evaluated result
\emph{term}, stored inline as a blob \src{sql:124}, not a hash into $\UStore$; caching a result does
not put it in the store.
\end{definition}

\begin{remark}[The key need not be a definition]\label{rem:watchkey}
$\dom(\UCache)$ identifies an expression, which need not be named or declared. The schema note says
the key ``is an unevaluated term reference \dots\ because the unevaluated term may not exist in the
codebase: it is not added merely by watching it without a name, e.g \texttt{> 2 + 3}''
\src{sql:141-143}. \Cref{prop:watchdom} is that sentence as a statement about the model. It is
observable in a running system: after evaluating one watch, the codebase's \texttt{hash} table has
interned the watch's hash and its \texttt{hash\_object} table has no row for it. Stored data carries
no notion of ``is a definition'' that the cache would have to respect, which is what lets a watch
expression in a scratch file be cached on the same footing as a top-level binding.
\end{remark}

\begin{theorem}[Cache soundness is stable under growth]\label{thm:wstable}
If $\UCache(r) = v$ is sound in $\UStore$ and $\UStore \subseteq \UStore'$, it is sound in
$\UStore'$.
\end{theorem}
\begin{proof}[Proof sketch]
Evaluating $r$ reads only entries reachable from $r$; all of them are in $\dom(\UStore)$ and, by
immutability, identical in $\UStore'$. The argument uses neither well-formedness nor anything about
the namespace, and is indifferent to whether the codomain is a hash or an inlined term.
\end{proof}

\begin{remark}[Why \cref{prop:watchdom} costs the cache nothing]\label{rem:watchdom-ok}
The proof above did not use that $r$ itself is in the store, only that everything $r$ reads is. A
cached key that names nothing, which \cref{prop:watchdom} shows is reachable, lies outside the
argument's hypotheses; it is not an edge case to exclude. Had the cache been specified as a map on
definitions, that proposition would have been a bug report.
\end{remark}

\begin{corollary}[No invalidation is needed]\label{cor:wnoinval}
An entry, once written, is valid in every future store. Consequently no invalidation logic is
required, and none exists: a search across the storage, codebase and CLI layers finds none. Two
removal paths exist and neither is invalidation: a debug-only wholesale \texttt{clear\-Watches},
reachable as \texttt{debug.clear-cache}, and a garbage collection of watches whose hash has no
object, which deletes from the \texttt{watch} table, not \texttt{watch\_result}, and whose only
caller is a one-time schema migration.
\end{corollary}

\begin{remark}[What ``never invalidated'' actually means]\label{rem:cachemeaning}
This is the claim Unison is best known for, and the informal version invites a weaker reading. It is
not an empirical observation that definitions rarely change, nor a policy traded against freshness.
It follows from two facts: the store is append-only, and the cache is keyed by a hash of the
expression that was evaluated. The argument mentions no names, types, or edit frequency.
\end{remark}

% =====================================================================
\section{Update}\label{sec:uupdate}

\subsection{Dependents}

\begin{definition}[Dependents]\label{def:dependents}
$\deps(h,i)$ is the set of references occurring in the $i$-th member of the component at $h$;
$\dependents$ is its converse, maintained as an index populated when a definition is saved. The index
can exist because a reference remains a node in the stored representation even though it hashes
transparently at position $0$ (\cref{prop:transparent}). Transparency is a property of the hash, not
of the representation.
\end{definition}

\begin{definition}[Update scope]\label{def:uscope}
\src{ops:1125-1126; update:163-179; branchtype:280-289} Given a set $Q$ of references being updated
and a namespace $\UNames$,
\[
  \tdeps(Q) \;=\; \bigl(\dependents^{+}(Q) \;\cap\; \mathcal{S}\bigr) \setminus F
\]
where $\dependents^{+}$ is the transitive closure excluding self-references, $F$ is what the incoming
file itself shadows, and $\mathcal{S}$ is the set of references named in $\UNames$ once $\lib$ has
been removed. Neither side condition is a filter on the dependents. $\lib$ is deleted from the
namespace before $\mathcal{S}$ is read off it. Unconflictedness is also not a filter: by
\cref{def:unconf} the implementation cannot take the unconflicted part of a namespace, so
$\mathcal{S}$ is well defined only because \cref{def:update} has already demanded
$\unconf{\UNames}{}$ of the whole non-$\lib$ namespace. A separate rule, not part of $\tdeps$, rejects
the update outright if the incoming file binds anything under $\lib$ \src{update:139-144}.
\end{definition}

\begin{remark}[Naming is what puts a definition in scope]\label{rem:scopenaming}
An unnamed definition is not a dependent to be updated, even if it references the target. Only naming
makes a stored object a participant in an edit. As in \cref{rem:watchkey}, the namespace carries
``definition'' status and stored data does not.

There is no type-based gating. Unison formerly classified each edit as
$\mathsf{Same} \mid \mathsf{Subtype} \mid \mathsf{Different}$ and propagated only the first two
automatically. The propagation machinery that consumed that classification was deleted in 2025; the
datatype survives, with no callers, because old history has to keep deserializing. The current design
computes the closure and lets the typechecker decide, reporting failure instead of predicting it.
\end{remark}

\subsection{Update, as rendering}

\begin{definition}[Pretty-print environment]\label{def:ppe}
\src{update:429-463} A $\ppe$ assigns each reference a name to render it by, under two
suffixification strategies: definitions in the file are suffixified by name, definitions in the
codebase by hash. The asymmetry is forced. Aliases sharing a suffix may be referred to by that suffix
in the codebase but not in a file, where the parser reports an ambiguity. The source's example: with
\texttt{one.foo} and \texttt{two.foo} both in the file, \texttt{hey = foo + foo} fails to parse.
\end{definition}

\begin{definition}[Update]\label{def:update}
\src{update:89, 219-300} Given a typechecked incoming file $F$ and a
configuration $\UConfig{\UStore}{\UNames}{\UCache}$ with $\unconf{\UNames}{}$:
\begin{enumerate}[nosep,label=\arabic*.]
  \item Let $Q$ be the references of the names $F$ redefines, and $D = \tdeps(Q)$.
  \item Render: $S = \render(\ppe, D)$, source text for every dependent.
  \item Parse and typecheck $F \cup S$ \emph{as one file}.
  \item If it typechecks, \textbf{store} each definition's component at the hash
        \cref{def:usave} already gave it, and rebind every affected name. This step is the
        \emph{commit}. By \cref{rem:saveident} the hashes are not computed here, only acted on.
  \item If it does not, nothing is ingested; \textbf{remove from $\UNames$ every name $F$ binds and
        every name in $D$}, and deliver $F \cup S$ to the user as editable source on a dedicated
        update branch.
\end{enumerate}
\end{definition}

\begin{remark}[How much step 5 removes]\label{rem:step5}
Step 5 removes more than the natural reading suggests. The implementation subtracts, by name,
everything the file binds together with every dependent it rendered, whatever the typechecker's
verdict on each. $F$'s own definitions do not survive either: since nothing is ingested, the name the
user was updating is left denoting nothing, not its old referent. This keeps the operation coherent,
since the namespace never holds a mixture of old and new. \Cref{ex:update-fail} works it through.
\end{remark}

\begin{remark}[The failure case is the interesting one]\label{rem:residual}
Informal descriptions of \texttt{update} rarely anticipate step 5. When propagation fails, the
operation neither commits broken definitions nor rolls back. The whole edit, what the user wrote and
everything it touched, is evicted from the namespace, and its source becomes the user's working
state. A rollback would restore the pre-edit namespace and leave the user's work only in whatever
buffer they typed it into; a commit would leave broken bindings. Step 5 makes the namespace smaller
and hands the difference to the user as code. What is handed back has to be readable, so this
depends on binder names surviving storage (\cref{rem:names}).

The same closure computation feeds \texttt{merge}, which puts conflicted definitions and their
dependents into a scratch file together for the same reason: whatever the user picks, everything
downstream must be re-typechecked against it.
\end{remark}

\begin{proposition}[Update propagates by exactly one strategy]\label{prop:follow-only}
The transition of \cref{def:update} admits no parameter controlling propagation: $\tdeps$ is fixed.
Every transitive dependent with an unconflicted name outside $\lib$ is rebuilt, and there is no way to
\emph{request} that some callers keep pointing at the old definition. One can be brought about anyway.
Give the definition a second name before updating. $\tdeps$ still selects every dependent and each is
still rebuilt, but rebuilding means rendering, and the renderer shows the dependency under one of its
names. The alias still denotes the old definition, so a dependent rendered under it parses back to
what it already said and the retarget is a no-op. The naming layer pins the caller, not an option.
\Cref{ex:merge-alias} exhibits it. Abstention is also available, by choosing a name nothing depends
on so that $\tdeps(Q)$ is empty. It is not a separate command either: at this commit \texttt{add} is
an alias for \texttt{update} \src{inputpat:882-896}, the two having been merged in 2025. A
per-definition choice does reappear at failure time, since whatever fails to typecheck comes back as
source for the user to edit.
\end{proposition}

\subsection{The obligation update rests on}

A stored definition refers to its dependencies by hash, so rebuilding one to call a new version means
changing a hash inside it. \Cref{def:update} does this without touching the stored tree. It prints
the definition, so the dependency becomes a name again, and parses it back in a namespace where that
name denotes the new version. The printer and parser perform the substitution, so the operation is
only as correct as the round trip between them.

\begin{proposition}[Update is correct only if printing inverts parsing]\label{prop:print-parse}
Let $d$ be a dependent not redefined by $F$, and let $d'$ be the definition obtained from step 4 for
$d$. The intended semantics of \cref{def:update} is that $d'$ differs from $d$ \emph{only} by having
its references to $Q$ replaced by their new targets. This holds iff, for every $d \in D$,
\[
  \uparse(\ppe,\ \render(\ppe, d)) \;=\; d
  \qquad\text{modulo exactly that replacement.}
\]
\end{proposition}
\begin{proof}[Proof sketch]
Step 3 reconstructs $d'$ from text alone. Any discrepancy the round trip introduces (a binder renamed
so as to change capture, an identifier resolving to a different alias, a different associativity) is
preserved into $d'$ and changes its hash for a reason unrelated to the update. Conversely, if the
round trip is exact except at the references being retargeted, $d'$ is by construction the intended
substitution.
\end{proof}

\begin{remark}[The property is understood precisely, and lives in a test]\label{rem:printparse}
The implementation's comment for this step reads: ``We are updating old references to new references
by rendering old references as names that are then parsed back to resolve to new references (the
world's weirdest implementation of AST substitution)'' \src{update:424-427}. The correctness property
that comment presupposes is not stated anywhere in the source. It is tested:
\texttt{unison-src/transcripts-round-trip/} renders a regression corpus, re-adds it to a fresh
namespace, and asserts that the namespace diff is empty (``this diff should be empty if the two
namespaces are equivalent. If it's nonempty, the diff will show us the hashes that differ.''). The
corpus file is named \texttt{reparses-with-same-hash.u}, and its sibling \texttt{reparses.u} holds the
weaker cases that need only parse.

The file names draw the distinction between ``reparses'' and ``reparses with the same hash'', so the
property is understood; it is maintained as a golden corpus instead of stated as an invariant. Given
the size of Unison's surface language that is a defensible choice. \Cref{prop:print-parse} makes
explicit what the corpus protects, and the difficulty concentrates in the pretty-print environment of
\cref{def:ppe}.
\end{remark}

% =====================================================================
\section{Merge}\label{sec:umerge}

\texttt{update} rebuilds a namespace after one person edits it. \texttt{merge} reconciles two
namespaces edited independently. Both render stored definitions to source and parse them back, so
both carry the obligation of \cref{prop:print-parse}. Merge also has to decide something
\texttt{update} never does: given that a definition differs on both branches, did a person change it,
or did something underneath it change? A second hash, computed from names, makes that decision.

\subsection{Three namespaces, one store}

A branch is modeled as a namespace. The causal history Unison attaches to one is out of scope
(\cref{sec:unotmodeled}), and so is computing two branches' common ancestor. Merge is a function of
three namespaces in the sense of \cref{def:unames} over one store: $\NO$ (the common ancestor), $\NA$
(the branch being merged \emph{into}) and $\NB$ (the branch being merged \emph{from}). The asymmetry
between the last two is retained throughout. Several steps prefer $\NA$, and \cref{prop:merge-asym} is
the consequence.

\begin{proposition}[The store never merges]\label{prop:merge-store}
Let $\UStore^{\mathsf{a}}$ and $\UStore^{\mathsf{b}}$ arise from independent sequences of ingests
(\cref{def:uingest}) starting from a common $\UStore^{\mathsf{o}}$. Then
$\UStore^{\mathsf{a}} \cup \UStore^{\mathsf{b}}$ is a store in the sense of \cref{def:ustore}, and
merge performs no reconciliation on it.
\end{proposition}
\begin{proof}
Well-definedness: if $h \in \dom(\UStore^{\mathsf{a}}) \cap \dom(\UStore^{\mathsf{b}})$ then both
components hash to $h$, so by \cref{thm:ualpha} they are the same component up to $\alpha$-equivalence
and cycle permutation, which is the identification $\comphash$ makes. No-dangling: each store is
closed under the references its components contain, and closure is preserved by union.
\end{proof}

\begin{remark}[What that proposition is and is not]\label{rem:merge-store}
It is the formal content of the claim, often made for content-addressed codebases, that two people
cannot conflict in the code. The store is append-only and keyed by content, so a union is always a
store and never needs a decision.

It is also close to vacuous. Nothing a user would call a merge conflict lives in $\UStore$. Conflicts
live in $\UNames$, which is not content-addressed: it is a relation over names people choose. The
rest of this section concerns $\UNames$, and none of it follows from \cref{prop:merge-store}.
\end{remark}

\begin{definition}[Merge preconditions]\label{def:mergepre}
\src{merge2:208-222; mblob:98-102} Merge demands three things before it computes anything.
\begin{enumerate}[nosep,label=(\roman*)]
  \item $\unconf{\NO}{}$, $\unconf{\NA}{}$ and $\unconf{\NB}{}$. By \cref{def:unconf} this is
        all-or-nothing on each of the three. A single conflicted name anywhere aborts the merge, and
        for $\NO$ the user cannot act on the diagnosis.
  \item Neither $\NA$ nor $\NB$ binds anything under $\lib$, the same restriction
        \cref{def:uscope} imposes on \texttt{update}.
  \item No \emph{conflicted alias}: no pair of names that denote one referent in $\NO$ and different
        ones on a branch. If \texttt{foo} and \texttt{bar} were aliases and one branch left them
        denoting different things, merge stops.
\end{enumerate}
Condition (iii) makes the rest of the section well posed. Everything below indexes by name, so a
definition must not be updatable at one of its names and not at another.
\end{definition}

\subsection{The syntactic hash}

\begin{definition}[Naming function]\label{def:namer}
\src{mdiffblob:412-422; names:359-369} $\namer : \Referent \rightharpoonup \Name$ assigns each
referent one \emph{fully qualified} name, drawn from $\NA$ if $\NA$ names it at all, otherwise from
$\NB$, otherwise from $\NO$. The bias is by referent, not by name: if $\NA$ has any name for $r$, none
of $\NB$'s names for $r$ are considered. Referents named nowhere are outside $\dom(\namer)$ and are
treated below as denoting themselves.
\end{definition}

\begin{definition}[Syntactic hash]\label{def:synhash}
\src{msynhash:237-246, 327-329, 386-390} For a resolved term $t$ (\cref{def:urterm}),
$\syn{\namer}{t}$ is computed by the recursion of \cref{def:uhash} with one substitution reversed: a
bound variable contributes its de Bruijn index, as before, and a reference $r$ contributes
$\namer(r)$ as \emph{text}, or its hash when $r \notin \dom(\namer)$. Binder names are discarded;
reference hashes are not consulted.
\end{definition}

\begin{remark}[The two hashes are duals, and only one is stored]\label{rem:syn-dual}
\Cref{def:uhash} erases names and keeps hashes. \Cref{def:synhash} erases hashes and keeps names.
They run over the same syntax with the same treatment of binders and differ in one clause. The
content hash is the definition's identity and is stored. The syntactic hash is computed during a
merge, used for one decision and discarded: it appears in no table, and no operation retrieves a
definition by it.

Two consequences follow. $\syn{\namer}{t}$ is a function of the namespace as well as the term, so the
same stored definition has different syntactic hashes in different namespaces. And equality of
syntactic hashes is equality of rendered form: a reader shown both under $\namer$ would see the same
code. \Cref{def:mergediff} is built on that reading.
\end{remark}

\begin{lemma}[Propagation is invisible to the syntactic hash]\label{lem:prop-silent}
Let $d'$ be obtained from $d$ by replacing each reference $r$ with some $r'$ satisfying
$\namer(r) = \namer(r')$. Then $\syn{\namer}{d} = \syn{\namer}{d'}$.
\end{lemma}
\begin{proof}
Immediate: by \cref{def:synhash} a reference contributes only through $\namer$, and the two terms
agree everywhere else.
\end{proof}

\begin{remark}[Why that lemma is the whole idea]\label{rem:prop-silent}
\Cref{def:update} rebuilds every transitive dependent of an edit, so after any edit many definitions
have new hashes that no human typed. By \cref{lem:prop-silent} those definitions are syntactically
unchanged, provided the edited definition kept its name so that its old and new versions have the
same $\namer$-image. The implementation's own illustration is this hypothesis: with
\texttt{foo = \#bar + 3} becoming \texttt{foo = \#bar2 + 3}, ``our pretty-print environment has mapped
both \texttt{\#bar} and \texttt{\#bar2} to the name \texttt{helper}, so each version of \texttt{foo}
would have the same syntactic hash. This indicates (to our merge algorithm) that this was an
auto-propagated update'' \src{msynhash:1-25}.
\end{remark}

\subsection{The diff, and what a conflict is}

\begin{definition}[Branch diff]\label{def:mergediff}
\src{mnarrow:29-32; mdiff:52-84} For a branch $P \in \{\NA, \NB\}$ and a name $x$, the change
$\delta_P(x)$ is
\[
\begin{array}{ll}
  \dadd & x \notin \dom(\NO),\ x \in \dom(P) \\
  \ddel & x \in \dom(\NO),\ x \notin \dom(P) \\
  \dnil & \NO(x) = P(x) \\
  \dprop & \NO(x) \neq P(x) \text{ and } \syn{\namer}{\NO(x)} = \syn{\namer}{P(x)} \\
  \dupd & \NO(x) \neq P(x) \text{ and } \syn{\namer}{\NO(x)} \neq \syn{\namer}{P(x)}
\end{array}
\]
Only $\dadd$, $\ddel$ and $\dupd$ enter the diff. $\dnil$ is discarded before any hashing, since
equal referents have equal syntactic hashes and the content hash decides the case alone. $\dprop$ is
partitioned off as an automatic propagation. It is held back from the diff, not discarded:
\cref{def:umerge} still applies the $\dprop$ bindings to the merged namespace, because a definition
\texttt{update} rebuilt is still the version that branch has. It is only spared from being put to
the user.
\end{definition}

\begin{remark}[Only the syntactic hash decides]\label{rem:syn-only}
\src{msynhashed:9-26} The $\dprop$/$\dupd$ distinction is the only place the algorithm consults a
definition's content hash, and even there it is consulted second. The type carrying a definition
through the diff pairs a referent with its syntactic hash, and its equality instance compares only the
syntactic hash: ``the \texttt{Eq} and \texttt{Ord} instances only compares syntactic hashes''. Once
two definitions' syntactic hashes differ, nothing downstream asks what they are. A content-addressed
system decides its merges by name.
\end{remark}

\begin{definition}[Combination]\label{def:mergecombine}
\src{mtwodiff:24-40; mcombine:51-67} A name $x$ is \emph{conflicted} iff the pair
$(\delta_{\NA}(x), \delta_{\NB}(x))$ is $(\dadd,\dadd)$ or $(\dupd,\dupd)$ and, in addition,
$\syn{\namer}{\NA(x)} \neq \syn{\namer}{\NB(x)}$. Every other pair resolves:
\begin{center}\small
\begin{tabular}{@{}ll@{}}
\toprule
Pair & Result \\
\midrule
one side only & that side's operation \\
$(\dadd,\dadd)$ or $(\dupd,\dupd)$, hashes equal & one operation, taking $\NA$'s referent \\
$(\ddel,\ddel)$ & delete \\
$(\ddel,\dupd)$ or $(\dupd,\ddel)$ & the update; the delete is discarded \\
\bottomrule
\end{tabular}
\end{center}
The merged namespace $\NM$ is $\NO$ with the resolved operations applied.
\end{definition}

\begin{remark}[A delete never wins]\label{rem:updel}
The source says the fourth row is a choice: ``These two are not a conflicts, perhaps only temporarily,
because it's easier to implement. We just ignore these deletes and keep the updates''
\src{mcombine:60-63}. It matters because Unison has no rename operation at the level of this diff:
moving \texttt{foo} to \texttt{bar} appears as $\ddel$ at \texttt{foo} and $\dadd$ at \texttt{bar}.
A rename racing an ordinary edit of the renamed definition therefore resolves by keeping both halves,
the new name from the rename and the old name from the edit. \Cref{ex:merge-dup} works that through.
\end{remark}

\begin{definition}[Merge]\label{def:umerge}
\src{mblob:186-243; merge2:346-380} Given $\NO$, $\NA$, $\NB$ satisfying \cref{def:mergepre}:
\begin{enumerate}[nosep,label=\arabic*.]
  \item Compute $\delta_{\NA}$, $\delta_{\NB}$ and the conflicted names (\cref{def:mergecombine}).
        Form $\NM$ from $\NO$ by applying the $\dprop$ bindings, then the resolved operations on top;
        conflicted names are removed, not bound.
  \item Call a referent \emph{superseded} if $\NO$ bound some name to it and $\NM$ no longer does.
        Let $D$ be every name $\NM$ binds whose definition reaches a superseded referent, together
        with the conflicted names. These must be rebuilt. An update supersedes (the name means
        something else now), a rename supersedes (the name is gone), and a definition all of whose
        dependencies still mean what they meant needs no attention.
  \item Render $D$ to source under a naming function $\rnamer$, producing one file.
  \item If no name is conflicted, parse and typecheck that file against $\NM$, from which the names
        in $D$ have been removed so that the re-ingest rebinds them, as \texttt{update} treats what
        it rebuilds. On success the merge is complete.
  \item Otherwise (a conflicted name, or a file that does not parse or typecheck), create a
        temporary branch holding $\NM$ and deliver the file to the user as editable source.
\end{enumerate}
$\rnamer$ is built like $\namer$ but suffixified, and prefers $\NA$'s names so that the result should
``look familiar to Alice (the one merging in Bob's changes)'' and avoid ``superfluous textual
conflicts that would arise from preferring Bob's names for Bob's code'' \src{mrender:93-104}.
\end{definition}

\begin{remark}[Two failures, one message]\label{rem:merge-twofail}
Step 5 covers two situations that the implementation reports identically. Either a name is
conflicted, in which case the file contains two definitions under that name and cannot parse by
construction; or no name is conflicted and the file still fails, because rendering and re-parsing did
not round-trip. \Cref{prop:merge-print-parse} concerns the second, and \cref{ex:merge-dangling}
exhibits it. The user sees the same sentence either way and must read the file to tell which happened.
\end{remark}

\subsection{What merge guarantees}

\begin{theorem}[A conflict requires two syntactic changes]\label{thm:noconflict}
If $\delta_{\NA}(x) \in \{\dnil, \dprop, \ddel\}$ then $x$ is not conflicted, whatever $\NB$ did.
Symmetrically for $\NB$.
\end{theorem}
\begin{proof}
$\dnil$ and $\dprop$ do not enter the diff, so $x$ is absent from $\NA$'s side of the combination and
falls in the ``one side only'' row of \cref{def:mergecombine}. $\ddel$ pairs only as $(\ddel,\ddel)$
or $(\ddel,\dupd)$, neither of which is a conflict.
\end{proof}

\begin{corollary}[Propagation alone never conflicts]\label{cor:noedit}
Suppose $\NA$'s binding at $x$ was not authored on $\NA$: either $\NA(x) = \NO(x)$, or $\NA(x)$ is
obtained from $\NO(x)$ by replacing each reference $r$ with an $r'$ satisfying
$\namer(r) = \namer(r')$. Then $x$ is not conflicted.
\end{corollary}
\begin{proof}
The first case gives $\delta_{\NA}(x) = \dnil$; the second gives $\dprop$ by
\cref{lem:prop-silent}. Apply \cref{thm:noconflict}.
\end{proof}

\begin{definition}[Authored; name-stable]\label{def:namestable}
\Cref{cor:noedit}'s hypothesis is about $\namer$, which no user thinks about. Restated as a condition
on what people did: fix a branch $P$ reached from $\NO$ by a sequence of gestures, each an
\texttt{update} (\cref{def:update}) or a \texttt{move}, \texttt{alias} or \texttt{delete}
(\cref{def:gestures}).
\begin{itemize}[nosep]
  \item A name $x$ is \textbf{authored on $P$} if some \texttt{update} in the sequence submitted a
        file binding $x$. Names that changed only because \texttt{update} rebuilt their dependents
        are not authored.
  \item $P$ is \textbf{name-stable} if no gesture is a \texttt{move}, an \texttt{alias} or a
        \texttt{delete}, and no \texttt{update} binds a name to a definition that already has one.
        The last clause is not redundant with \texttt{alias}: by \cref{thm:ualpha}, writing out a
        definition that already exists produces the same referent, so a duplicate is an alias arrived
        at without asking for one.
\end{itemize}
\end{definition}

\begin{corollary}[Nobody is asked about code nobody wrote, if nothing is renamed]\label{cor:noedit-user}
If $\NA$ and $\NB$ are both name-stable (\cref{def:namestable}), then no name authored on neither
branch is conflicted.
\end{corollary}
\begin{proof}
Let $x$ be authored on neither. If $x \notin \dom(\NO)$ then, since neither branch moves or aliases,
$x$ is bound on a branch only if authored there, so it appears in no diff. Otherwise $x \in \dom(\NO)$,
and by \cref{thm:noconflict} and \cref{cor:noedit} it is enough to show $\namer(r) = \namer(r')$ for
each reference $r$ of $\NO(x)$ that \cref{def:update} retargeted to $r'$.

Let $y$ be the name at which that dependency sits, so $\NO(y) = r$ and $\NA(y) = r'$. Since $\NA$ is
name-stable, $y$ is its only name for $r'$, so $\namer(r') = y$. $\NA$ has no name for $r$, since $y$
now denotes $r'$ and no alias or duplicate gave $r$ another. So $\namer(r)$ falls through to $\NB$,
which by its own name-stability names $r$ at $y$ if it left the dependency alone, or nowhere if it
updated it. In the latter case $\namer(r)$ falls through again to $\NO$, where $r$ is named $y$.
Either way $\namer(r) = y = \namer(r')$.
\end{proof}

\begin{remark}[Which is to say: only a change of names does this]\label{rem:onlyrenames}
Contrapositively, the only user action that can put someone in front of a conflict in a definition
they never opened is one that changes which names a definition has. Adding, updating and deleting
definitions are safe at any depth and in any combination. The gestures of \cref{def:gestures} are the
exception, and they are the one class of gesture that costs nothing anywhere else in the system.

Three details of the hypothesis matter.

\emph{It binds both branches, not only the one that abstained.} If $\NA$ updates a dependency in place
and $\NB$ renames it, $\namer$ takes $\NB$'s new name for the old version and $\NA$'s unchanged name
for the new one, and $\NA$'s untouched dependents are classified as authored. $\NA$ can be fully
disciplined and still be asked about code it never wrote.

\emph{An alias suffices; a \texttt{move} is not needed}, which is why \cref{def:namestable} excludes
aliases. $\namer$ picks among a referent's names by a fixed priority, relative before absolute and
then fewest segments \src{ppenames:141-151}. Giving an updated definition a second, higher-priority
name changes what it renders as, while the superseded version keeps the original name from the other
branch. \Cref{ex:merge-alias} shows this, and also the reverse case: an alias created before the
update pins the dependents instead of conflicting them.

\emph{Unconflicted is not unseen.} Step 2 of \cref{def:umerge} pulls the dependents of every conflict
into the same file. \Cref{cor:noedit-user} rules out only being asked to choose.
\end{remark}

\begin{proposition}[Renaming defeats the side condition]\label{prop:rename-conflict}
There are configurations in which no definition at $x$ was authored on either branch beyond one
party's ordinary edit, and $x$ is nonetheless conflicted. It suffices that $\NA$ update a definition
$y$ on which $x$ depends and change $y$'s name in the same operation, and that $\NB$ edit $x$.
\end{proposition}
\begin{proof}
Let $\NO$ bind $y \mapsto r$ and $x \mapsto d$ with $r$ occurring in $d$. Let $\NA$ bind
$y' \mapsto r'$ and $x \mapsto d'$, where $d'$ is $d$ with $r$ replaced by $r'$, the propagation
\cref{def:update} performs. Since $\NA$ has no name for $r$ and $\NB$ does, $\namer(r) = y$; since
$\NA$ names $r'$, $\namer(r') = y'$. So $\syn{\namer}{d} \neq \syn{\namer}{d'}$ and
$\delta_{\NA}(x) = \dupd$. If $\NB$ edits $x$ then $\delta_{\NB}(x) = \dupd$ with a different
syntactic hash, and $x$ is conflicted by \cref{def:mergecombine}.
\end{proof}

\begin{remark}[Two gestures, one intention]\label{rem:rename-gesture}
The hypothesis is not exotic. Renaming is \texttt{move} and editing is \texttt{update}, two commands,
but ``this is badly named and also wrong'' is one thought, discharged by running them back to back. A
\texttt{move} with no \texttt{update} does not conflict: $\namer$ follows the referent, so the
ancestor's version of $x$ renders under the new name too and the two syntactic hashes agree. It fails
differently. \Crefrange{ex:merge-rename}{ex:merge-dup} are the three cases.
\end{remark}

\begin{corollary}[The side condition is computed, and discarded]\label{cor:unchecked}
\src{mrename:31-96; mdiffblob:171-177} Merge builds, for each definition, the set of names it gained
and the set it lost between $\NO$ and each branch, and narrows that to the one-to-one cases, the
renames. At \texttt{db60ce2} the result reaches a debug printer and nothing else: no part of the diff,
the combination, $\NM$ or the rendered file consults it.
\end{corollary}

\begin{remark}[Compare \cref{cor:detect}]\label{rem:two-side-conditions}
Two conditional guarantees are handled oppositely. \Cref{thm:ordering} fails when two members of a
cycle are indistinguishable; this is decided at ingest and the implementation aborts
(\cref{cor:detect}). \Cref{cor:noedit} fails when a name moves; this is computed and discarded, and
merge proceeds either way, reporting a conflict it cannot tell from a real one.

They are not equally tractable. \Cref{cor:unchecked}'s detection groups definitions by syntactic
hash, so it finds a name that moved while the definition stayed put, the case
\cref{rem:rename-gesture} says does not conflict. The case that does conflict is a rename plus an
edit, where the syntactic hash moves too and the two versions are never grouped. The existing
machinery addresses the harmless case; the harmful one needs a signal a diff of two namespaces does
not carry.
\end{remark}

\begin{remark}[The one place the condition is repaired]\label{rem:libcanon}
\src{mdiffblob:463-505} Before $\namer$ is built, names under $\lib$ have a trailing mangled version
number stripped, so that \texttt{base\_1\_0\_0} and \texttt{base\_2\_0\_0} become one name. The source
gives the reason, ``to make it more likely that we classify things that are \emph{probably} true
propagated updates as such'', and calls it a hack.

This restores \cref{cor:noedit}'s hypothesis by hand. Upgrading a library renames everything in it,
so without the stripping every definition mentioning the library classifies as authored and every
concurrent edit conflicts. Merge forces the two names equal instead of tracking the rename, which
works for the one renaming scheme whose shape is known in advance.
\end{remark}

\begin{proposition}[Merge is not commutative]\label{prop:merge-asym}
Merging $\NB$ into $\NA$ and merging $\NA$ into $\NB$ can yield different outcomes, including one
succeeding where the other fails, with no conflict reported in either direction.
\end{proposition}
\begin{proof}[Proof sketch]
Two mechanisms suffice independently. $\namer$ and $\rnamer$ both prefer $\NA$'s names by referent
(\cref{def:namer}), so both the classification of \cref{def:mergediff} and the text produced by
\cref{def:umerge} depend on which branch is the target. And the second row of
\cref{def:mergecombine} takes $\NA$'s referent when two syntactic hashes agree. \Cref{ex:merge-dangling}
exhibits a pair of branches for which one direction is a clean merge and the other is a failure.
\end{proof}

\begin{proposition}[Merge's round trip spans two namespaces]\label{prop:merge-print-parse}
Step 4 of \cref{def:umerge} requires, for every definition $d$ it renders,
\[
  \uparse(\NM,\ \render(\rnamer, d)) \;=\; d \qquad\text{modulo the intended retargeting,}
\]
which is \cref{prop:print-parse} with two differences. The environment that renders is not the
environment that parses: $\rnamer$ prefers $\NA$'s names, while $\NM$ has applied both branches'
deletions and additions. And nothing requires that a name $\rnamer$ emits still exists in $\NM$.
\end{proposition}
\begin{remark}[The gap is reachable]\label{rem:merge-gap}
If $\NB$ renames a definition and $\NA$ does not, $\NM$ holds $\NB$'s new name and not the old one,
since a rename is an unconflicted delete plus an unconflicted add and both halves apply. But $\rnamer$
prefers $\NA$'s names, and $\NA$ still uses the old one. The file is written with a name the
namespace it will be parsed against no longer has, and step 4 fails (\cref{ex:merge-dangling}).
Reversing the roles removes the mismatch, which is \cref{prop:merge-asym}.
\end{remark}

\begin{remark}[The silent resolution, and what it hides]\label{rem:merge-soundness}
\Cref{thm:noconflict} and its corollary bound spurious conflicts. The dual question is whether every
real disagreement is reported. The case to examine is the second row of \cref{def:mergecombine}: when
both branches update $x$ to definitions with equal syntactic hashes but different referents, merge
takes $\NA$'s and says nothing.

The case is reachable. Let both branches update a shared dependency $y$ differently. Each rebuilds
$x$ by \cref{def:update}, and both rebuilds render as the same code, because $\namer$ gives each
branch's version of $y$ the same name. So $x$ has two versions and one syntactic hash, and is resolved
silently.

Nothing is hidden, because the collapse cannot go all the way down.
$\syn{\namer}{\NA(x)} = \syn{\namer}{\NB(x)}$ says the two render alike; if their referents differ,
some dependency differs, and the same question recurs at that dependency's name. The recursion
terminates: \cref{def:synhash} records literals, structure and types verbatim, and a referent $\namer$
does not name contributes its hash, so a difference invisible at every level does not exist. It
surfaces at some name, as a conflict if the two versions render differently there, which is where the
recursion stops. In the two-definition case it is reported at $y$ while $x$ resolves silently, and
$x$ is not what the parties disagree about.

We state this as an argument, not a theorem. What is checked is the witness: the configuration above,
with the conflict landing on the dependency and not on the definition resolved for it.
\end{remark}

% =====================================================================
\section{Worked examples}\label{sec:examples}

Each example is a derivation: every line applies an operation defined above, with the definition
named on the right. Symbolic hashes are written \texttt{\#a}, \texttt{\#b}; $\ulcorner t \urcorner$
abbreviates $\uhash{\HEnv}{t}$ where the environment is clear; and $r_{+}$ abbreviates
$\UBuiltin{\texttt{Nat.+}}$, a builtin, which therefore never acquires a hash of its own.

\subsection{Ingesting one definition}\label{ex:ingest}

\[
\begin{array}{r@{\;}l@{\quad}l}
  e &= \ulam{n}{\uapp{\uapp{\uvar{\texttt{Nat.+}}}{\uvar{n}}}{\unat{1}}}
    & \just{source term, \cref{def:uterm}} \\[3pt]
  \Res{\{\texttt{inc}\}}{\UNames}(e)
    &= \ulam{n}{\uapp{\uapp{\uref{r_{+}}}{\uvar{n}}}{\unat{1}}}
    & \just{resolution, \cref{def:uresolution}} \\[3pt]
  {} &\utypes{\emptyset}{\Res{\{\texttt{inc}\}}{\UNames}(e)}{\uarr{\TNat}{\TNat}}
    & \just{typing, \cref{def:utyping}} \\[3pt]
  t &= \uann{\Res{\{\texttt{inc}\}}{\UNames}(e)}{\uarr{\TNat}{\TNat}}
    & \just{type incorporation, \cref{def:typeinhash}} \\[3pt]
  \bg &= \{\,\texttt{inc} \mapsto t\,\}
    & \just{binding group, \cref{def:bindgroup}} \\[3pt]
  C &= \langle(\texttt{inc}, t)\rangle,\quad k = 1
    & \just{one SCC, \cref{def:uingest}} \\[6pt]
  \phash{\texttt{inc}} &= \uhash{\mathsf{L}(\{\texttt{inc}\})}{t}
    & \just{step 1, \cref{def:comphash}} \\[3pt]
  \HEnv &= \mathsf{R}(\texttt{inc}),\qquad h_1 = \uhash{\HEnv}{t}
    & \just{step 2} \\[3pt]
  \nhash{\texttt{inc}} &= \mathsf{acc}[\mathsf{Tag}\,1,\ \mathsf{Hashed}\,h_1,\
                          \mathsf{Hashed}\,\uhash{\HEnv}{\uvar{\texttt{inc}}}]
    & \just{step 3} \\[3pt]
  \texttt{\#a} &= \comphash(C) = \mathsf{acc}[\mathsf{Hashed}\,\nhash{\texttt{inc}}]
    & \just{step 3} \\[6pt]
  \UStore' &= \upd{\UStore}{\texttt{\#a}}{\langle t \rangle}
    & \just{store, \cref{def:ustore}} \\[3pt]
  \UNames' &= \UNames \cup \{(\texttt{inc},\ \URefT{\UDer{\texttt{\#a}}{0}})\}
    & \just{naming, \cref{def:unames}}
\end{array}
\]

The name is attached last, to a definition that was complete without it. Step 3 applies two
accumulations, so $\comphash(C) \neq h_1$ even at $k = 1$; \cref{ex:transparency} turns on this.

\subsection{An alias, and what transparency gives}\label{ex:transparency}

\[
\begin{array}{r@{\;}l@{\quad}l}
  e_x &= \uapp{\uapp{\uvar{\texttt{Nat.+}}}{\unat{1}}}{\unat{1}},
  \qquad e_y = \uvar{\texttt{x}}
    & \just{source, \cref{def:uterm}} \\[3pt]
  \Res{\{\texttt{x},\texttt{y}\}}{\UNames}(e_y) &= \uvar{\texttt{x}}
    & \just{unresolved: $\texttt{x} \in B$, \cref{def:uresolution}} \\[3pt]
  \bg &= \{\,\texttt{x} \mapsto X,\ \ \texttt{y} \mapsto \uann{\uvar{\texttt{x}}}{\TNat}\,\}
    & \just{\cref{def:bindgroup}} \\[3pt]
  X &= \uann{(\uapp{\uapp{\uref{r_{+}}}{\unat 1}}{\unat 1})}{\TNat} & \\[3pt]
  C_1 &= \langle(\texttt{x}, X)\rangle,\qquad C_2 = \langle(\texttt{y}, \cdot)\rangle
    & \just{two SCCs, \cref{def:uingest}} \\[6pt]
  \texttt{\#x} &= \comphash(C_1)
    & \just{as in \cref{ex:ingest}} \\[3pt]
  Y &= \uann{\uref{\UDer{\texttt{\#x}}{0}}}{\TNat}
    & \just{substitution before $C_2$, \cref{def:uingest}} \\[6pt]
  \uhashc{\uref{\UDer{\texttt{\#x}}{0}}} &= \texttt{\#x}
    & \just{\cref{prop:transparent}} \\[3pt]
  \uhashc{Y} &= \mathsf{acc}[\ldots,\ \mathsf{Hashed}\,\texttt{\#x},\ \ldots]
    & \just{via \cref{prop:transparent}} \\[3pt]
  \uhashc{X} &= \mathsf{acc}[\ldots,\ \mathsf{Hashed}\,\uhashc{B_x},\ \ldots]
    & \just{same position} \\[3pt]
  B_x &= \uapp{\uapp{\uref{r_{+}}}{\unat 1}}{\unat 1} &
\end{array}
\]

\begin{remark}[The usual summary of transparency is too strong]\label{rem:transparency-limit}
The source comment illustrates transparency with this pair, saying that \texttt{x = 1 + 1} and
\texttt{y = x} ``hash the same''. Comparing the last two lines, that would require
\[
  \uhashc{\uapp{\uapp{\uref{r_{+}}}{\unat 1}}{\unat 1}} \;=\; \texttt{\#x} ,
\]
the referent's \emph{term} hash to equal its \emph{component} hash. But \texttt{\#x} is built from
$\uhashc{X}$ by the two accumulations at the end of \cref{ex:ingest}, so the equation asks the hash
function for a fixed point of a particular shape. Transparency passes along the component hash,
inlining would contribute the term hash, and \cref{def:comphash} puts two layers between them. The
two definitions receive different hashes, and inlining a reference into a larger term does not
preserve that term's hash either.

What holds is \cref{prop:transparent}: a position-$0$ reference is indistinguishable from any term
hashing to \texttt{\#x}. The comment's own caveat is this condition: ``so long as
\texttt{Reference.Derived} ctors are created using the same hashing function as is used here''.
\texttt{hashClosedTerm} satisfies it, building $\UDer{\uhashc{t}}{0}$ directly from a term's hash
\src{term:124-125}, and references made that way are interchangeable with their referents. The
component construction does not.

This is the reading we are least confident of, and it bears on a frequently repeated description of
the system. If the definition-level claim does hold, we have misread the relationship between
$\comphash$ and the member hashes.
\end{remark}

\subsection{A recursive group}\label{ex:cycle}

Write $\langle j \rangle$ for the hash of de Bruijn index $j$, and $V = \{\texttt{cd},\texttt{cu}\}$.

\[
\begin{array}{r@{\;}l@{\quad}l}
  e_{\texttt{cd}} &= \ulam{n}{\uapp{\uvar{\texttt{cu}}}{\uvar{n}}},
  \qquad
  e_{\texttt{cu}} = \ulam{n}{\uapp{\uvar{\texttt{cd}}}{(\uapp{\uapp{\uvar{\texttt{Nat.+}}}{\uvar n}}{\unat 1})}}
    & \just{source} \\[3pt]
  \Res{V}{\UNames}(e_{\texttt{cd}}) &= e_{\texttt{cd}}
    & \just{both names in $B$, \cref{def:uresolution}} \\[3pt]
  C &= \langle(\texttt{cd}, \cdot), (\texttt{cu}, \cdot)\rangle,\quad k = 2
    & \just{one SCC, \cref{def:uingest}} \\[6pt]
  \phash{\texttt{cd}} &= \uhash{\mathsf{L}(V)}{\ulam{n}{\uapp{\langle 1\rangle}{\langle 0\rangle}}}
    & \just{step 1, \cref{def:comphash}} \\[3pt]
  \phash{\texttt{cu}} &= \uhash{\mathsf{L}(V)}{\ulam{n}{\uapp{\langle 1\rangle}{(\uapp{\uapp{\uref{r_{+}}}{\langle 0\rangle}}{\unat 1})}}}
    & \just{step 1} \\[3pt]
  \phash{\texttt{cd}} &\neq \phash{\texttt{cu}}
    & \just{the bodies differ} \\[3pt]
  \HEnv &= \mathsf{R}(\texttt{cd})\,\mathsf{R}(\texttt{cu})
    & \just{step 2, sorted by $\phash{\cdot}$} \\[3pt]
  \sigma &= \text{sort by } \nhash{\cdot}
    & \just{step 3} \\[3pt]
  \texttt{\#c} &= \comphash(C)
    & \just{step 3}
\end{array}
\]

Neither name enters $\phash{\cdot}$, so structure alone decides the ordering, as
\cref{thm:ordering}(a) requires. The second sort $\sigma$ is by a hash of the first sort's index and
need not agree with it (\cref{rem:twoorders}), so \texttt{cu} may take position $0$. Whichever member
does is referenced transparently and the other opaquely (\cref{prop:transparent}). The user wrote the
two symmetrically and cannot observe the asymmetry.

\subsection{When the canonical order is ambiguous}\label{ex:ambiguous}

The two behaviours that \cref{thm:ordering}(b1) and (b2) separate have been computed as well as
derived; see \cref{sec:paths}.

Change the two so they differ only in which of the other they call:

\[
\begin{array}{r@{\;}l@{\quad}l}
  e_{\texttt{ping}} &= \ulam{n}{\uapp{\uvar{\texttt{pong}}}{\uvar{n}}},
  \qquad
  e_{\texttt{pong}} = \ulam{n}{\uapp{\uvar{\texttt{ping}}}{\uvar{n}}}
    & \just{source} \\[3pt]
  \phash{\texttt{ping}} &= \uhash{\mathsf{L}(V)}{\ulam{n}{\uapp{\langle 1\rangle}{\langle 0\rangle}}}
    & \just{step 1, \cref{def:comphash}} \\[3pt]
  \phash{\texttt{pong}} &= \uhash{\mathsf{L}(V)}{\ulam{n}{\uapp{\langle 1\rangle}{\langle 0\rangle}}}
    & \just{step 1} \\[3pt]
  \phash{\texttt{ping}} &= \phash{\texttt{pong}}
    & \just{$\phash{\cdot}$ not injective} \\[3pt]
  \HEnv &= \mathsf{R}(\texttt{ping})\,\mathsf{R}(\texttt{pong})
    & \just{stable sort $\Rightarrow$ name order}
\end{array}
\]

In $\mathsf{L}(V)$ both mutual references collapse to one index, so the two bodies become the same
term and the sort cannot separate them. It is stable, and the list arrived in name order, so
\texttt{ping} precedes \texttt{pong} only because of the alphabet.

Renaming \texttt{ping} to \texttt{zing} swaps the binding order, but the obvious conclusion is wrong.
The two bodies are $\ulam{n}{\uapp{\langle 1\rangle}{\langle 0\rangle}}$ under either listing, so the
swap does not change the multiset $\{h_1,h_2\}$. It exchanges which name owns which index, hence
which $\nhash{\cdot}$, and \cref{def:comphash} accumulates those sorted. So \texttt{\#c} is unchanged,
and what moves is \textbf{which member holds position $0$}: which one is referenced transparently and
which opaquely. That is \cref{thm:ordering}(b1), the consequence that always holds.

For (b2), where $\comphash$ itself moves, the exchange has to change a member hash and not merely
relabel it. Two bodies that mention the cycle asymmetrically do that:
\[
  \texttt{foo} = \uapp{\uapp{\uvar{\texttt{Nat.+}}}{(\uapp{\uvar{\texttt{foo}}}{\unat 0})}}{(\uapp{\uvar{\texttt{bar}}}{\unat 0})},
  \qquad
  \texttt{bar} = \uapp{\uapp{\uvar{\texttt{Nat.+}}}{(\uapp{\uvar{\texttt{foo}}}{\unat 0})}}{(\uapp{\uvar{\texttt{bar}}}{\unat 0})}
\]
are tied under $\mathsf{L}(V)$ but, in $\HEnv$, hash to a term mentioning $\langle 1 \rangle$ before
$\langle 0 \rangle$ or the reverse depending on which name took the lower index. Renaming
\texttt{bar} to \texttt{zar} therefore changes $h_1$ and $h_2$, hence the shared token prefix
$\overline{\tau}$, hence \texttt{\#c}.

Either way the implementation detects the tie (\cref{cor:detect}) and aborts instead of storing a
name-dependent identity. It suggests adding something like
$\texttt{\_} = \utext{this is the ping definition}$, which makes $\phash{\cdot}$ injective again. It
aborts on the tie, not on (b2), declining the whole class instead of telling the harmless case from
the harmful one; \cref{rem:notpartial} argues this is right.

\subsection{An update that succeeds}\label{ex:update-ok}

Starting state, from two earlier ingests:
$\texttt{scale} \mapsto \UDer{\texttt{\#s}}{0}$ and $\texttt{double} \mapsto \UDer{\texttt{\#d}}{0}$,
where \texttt{double}'s stored body is $\ulam{n}{\uapp{\uref{\UDer{\texttt{\#s}}{0}}}{\uvar n}}$. The
user submits a new \texttt{scale} of the same type, $\uarr{\TNat}{\TNat}$.

\[
\begin{array}{r@{\;}l@{\quad}l}
  Q &= \{\UDer{\texttt{\#s}}{0}\}
    & \just{what $F$ redefines, \cref{def:update}} \\[3pt]
  \tdeps(Q) &= \{\UDer{\texttt{\#d}}{0}\}
    & \just{\cref{def:uscope}: named, unconflicted, not in $\lib$} \\[3pt]
  S &= \render(\ppe, \tdeps(Q))
    & \just{\cref{def:ppe}} \\[3pt]
  {} &= ``\texttt{double = }\ulam{n}{\uapp{\uvar{\texttt{scale}}}{\uvar n}}\text{''}
    & \just{a \emph{name}, not a hash} \\[3pt]
  \uparse(\ppe, S) &\rightsquigarrow \ulam{n}{\uapp{\uvar{\texttt{scale}}}{\uvar n}}
    & \just{a \emph{variable} again, \cref{def:uterm}} \\[3pt]
  \Res{\emptyset}{\UNames}(\cdot) &= \ulam{n}{\uapp{\uref{\UDer{\texttt{\#s}'}{0}}}{\uvar n}}
    & \just{\emph{here} the reference is retargeted} \\[3pt]
  {} &\utypes{\emptyset}{F \cup S}{\ } \text{ succeeds}
    & \just{\cref{def:utyping}} \\[3pt]
  \texttt{\#s}',\ \texttt{\#d}' &= \hashcomps(\bg_{F \cup S})
    & \just{\cref{def:uingest}} \\[3pt]
  \UNames'' &= \UNames[\texttt{scale} \mapsto \UDer{\texttt{\#s}'}{0}]
    & \just{\cref{def:update}, step 4} \\[3pt]
  \UNames''' &= \UNames''[\texttt{double} \mapsto \UDer{\texttt{\#d}'}{0}] &
\end{array}
\]

The fifth line carries the mechanism: retargeting happens because a printed name resolves against the
new namespace. \texttt{double}'s source text is unchanged, character for character, and its hash
changed anyway. \texttt{\#d} is still in $\UStore$, still referring to \texttt{\#s}, so any cached
result for it remains valid (\cref{cor:wnoinval}).

\subsection{An update that fails}\label{ex:update-fail}

Same starting state; this time the new \texttt{scale} is
$\ulam{s}{\uapp{\uvar{\texttt{Text.size}}}{\uvar s}}$, of type $\uarr{\TText}{\TNat}$.

\[
\begin{array}{r@{\;}l@{\quad}l}
  Q,\ \tdeps(Q),\ S &\text{ as before}
    & \just{\cref{def:uscope,def:ppe}} \\[3pt]
  {} &\utypes{\emptyset}{F \cup S}{\ } \text{ fails at \texttt{double}}
    & \just{$n : \TNat$ against $\uarr{\TText}{\TNat}$} \\[3pt]
  \UStore' &= \UStore
    & \just{nothing is ingested, \cref{rem:step5}} \\[3pt]
  \UNames'' &= \UNames \setminus \{(\texttt{scale}, \cdot),\ (\texttt{double}, \cdot)\}
    & \just{\cref{def:update}, step 5} \\[3pt]
  \text{delivered} &= F \cup S
    & \just{step 5}
\end{array}
\]

\texttt{scale} is neither rebound to a new hash nor left denoting its old one. It is among the names
$F$ binds, so step 5 removes it, and since nothing was ingested there is no \texttt{\#s}$'$ for it to
denote. After a failed update the namespace contains neither name.

The namespace holds no name denoting a definition that fails to typecheck, and no half-applied edit.
\texttt{\#s} and \texttt{\#d} are neither deleted nor altered; they are no longer reachable by name,
which is the only kind of deletion the store supports. The unfinished work is a file containing both
the user's new \texttt{scale} and the \texttt{double} it broke (\cref{rem:residual}). The third line
rests on \cref{prop:print-parse}: $\render$ must produce something $\uparse$ turns back into
\texttt{double}, differing only in what \texttt{scale} resolves to.

\subsection{Three merges from one ancestor}\label{ex:merge-setup}

The next four examples share a common ancestor with two definitions, a formatter and its one caller:
\[
  \NO = \{\,\texttt{fmt} \mapsto \texttt{\#f},\quad
           \texttt{show} \mapsto \texttt{\#c}\,\}
  \qquad
  \UStore(\texttt{\#c})_0 = \ulam{n}{\uapp{\uref{\texttt{\#f}}}{\uvar n}}
\]
with $\texttt{fmt} : \uarr{\TNat}{\TText}$. Every component here is a singleton, so throughout these
examples we write $\texttt{\#f}$ for $\UDer{\texttt{\#f}}{0}$. In each, one branch is the merge
target $\NA$ and the other the source $\NB$. Nobody edits a definition twice, and every branch is the
result of a single \texttt{update}, except in \cref{ex:merge-alias}, which needs one \texttt{alias}.
The examples differ only in who does what.

\subsection{A rename manufactures a conflict}\label{ex:merge-rename}

$\NA$ rewrites the formatter and renames it in the same edit; \cref{def:update} propagates to
\texttt{show}, which $\NA$ never opens. $\NB$ edits \texttt{show} and leaves the formatter alone.
\[
  \NA = \{\texttt{price} \mapsto \texttt{\#f}',\ \texttt{show} \mapsto \texttt{\#c}'\}
  \qquad
  \NB = \{\texttt{fmt} \mapsto \texttt{\#f},\ \texttt{show} \mapsto \texttt{\#c}''\}
\]
\[
\begin{array}{r@{\;}l@{\quad}l}
  \namer(\texttt{\#f}') &= \texttt{price}
    & \just{$\NA$ names it, \cref{def:namer}} \\[3pt]
  \namer(\texttt{\#f}) &= \texttt{fmt}
    & \just{$\NA$ does not; $\NB$ does} \\[3pt]
  \syn{\namer}{\NO(\texttt{show})} &= \syn{\namer}{\ulam{n}{\uapp{\uvar{\texttt{fmt}}}{\uvar n}}} & \\[3pt]
  \syn{\namer}{\NA(\texttt{show})} &= \syn{\namer}{\ulam{n}{\uapp{\uvar{\texttt{price}}}{\uvar n}}}
    & \just{same code, one identifier} \\[3pt]
  \delta_{\NA}(\texttt{show}) &= \dupd
    & \just{\cref{def:mergediff}: the two differ} \\[3pt]
  \delta_{\NB}(\texttt{show}) &= \dupd & \\[3pt]
  \texttt{show} &\text{ is conflicted}
    & \just{\cref{def:mergecombine}}
\end{array}
\]
The failure is in the fourth line. $\NA$'s \texttt{show} is the ancestor's own definition with one
identifier substituted by propagation, and it is presented to the user as one side of a conflict to be
resolved against $\NB$'s. Moreover, $\delta_{\NA}(\texttt{fmt}) = \ddel$ resolves unopposed, so $\NM$
has no \texttt{fmt}, yet $\NB$'s half of the conflict, which the user must now edit, refers to
\texttt{fmt}. Whichever side the user keeps, the rename has to be applied by hand, which is the work
\cref{prop:merge-store} was supposed to have removed.

\subsection{A rename alone: no conflict, and a merge that still fails}\label{ex:merge-dangling}

Now the rename carries no edit. $\NB$ renames the formatter and stops; $\NA$ edits \texttt{show}.
\[
  \NA = \{\texttt{fmt} \mapsto \texttt{\#f},\ \texttt{show} \mapsto \texttt{\#c}'\}
  \qquad
  \NB = \{\texttt{price} \mapsto \texttt{\#f},\ \texttt{show} \mapsto \texttt{\#c}\}
\]
\[
\begin{array}{r@{\;}l@{\quad}l}
  \delta_{\NA} &= \{\texttt{show} \mapsto \dupd\}
    & \just{$\texttt{fmt}$ unchanged, so $\dnil$} \\[3pt]
  \delta_{\NB} &= \{\texttt{fmt} \mapsto \ddel,\ \texttt{price} \mapsto \dadd\}
    & \just{a rename is a delete and an add} \\[3pt]
  \text{conflicted} &= \emptyset
    & \just{no name appears on both sides} \\[3pt]
  \NM &= \{\texttt{price} \mapsto \texttt{\#f},\ \texttt{show} \mapsto \texttt{\#c}'\} & \\[3pt]
  \rnamer(\texttt{\#f}) &= \texttt{fmt}
    & \just{$\rnamer$ prefers $\NA$, \cref{def:umerge}} \\[3pt]
  \render(\rnamer, \NM(\texttt{show})) &= ``\texttt{show = }\ulam{n}{\uapp{\uvar{\texttt{fmt}}}{\uvar n}}\text{''} & \\[3pt]
  \uparse(\NM, \cdot) &\text{ fails}
    & \just{$\texttt{fmt} \notin \dom(\NM)$}
\end{array}
\]
No name is conflicted and the merge fails anyway, at step 4 of \cref{def:umerge}, not step 1. The
file handed back contains a single \texttt{show}, not the two-sided form a conflict produces, written
with a name the branch it was written for does not have.

Exchanging the roles fixes it. With $\NB$ as the target, $\rnamer$ prefers \texttt{price}, the render
agrees with $\NM$, and the merge completes cleanly with \texttt{show} rebound. The same two branches,
merged in the two directions, give a failure and a success (\cref{prop:merge-asym}).

\subsection{A rename raced by an edit, silently undone}\label{ex:merge-dup}

$\NA$ renames the formatter and stops. $\NB$ edits the formatter's body, and \cref{def:update}
propagates to \texttt{show}.
\[
  \NA = \{\texttt{price} \mapsto \texttt{\#f},\ \texttt{show} \mapsto \texttt{\#c}\}
  \qquad
  \NB = \{\texttt{fmt} \mapsto \texttt{\#f}',\ \texttt{show} \mapsto \texttt{\#c}'\}
\]
\[
\begin{array}{r@{\;}l@{\quad}l}
  \delta_{\NA} &= \{\texttt{fmt} \mapsto \ddel,\ \texttt{price} \mapsto \dadd\} & \\[3pt]
  \delta_{\NB}(\texttt{fmt}) &= \dupd & \\[3pt]
  \delta_{\NB}(\texttt{show}) &= \dupd
    & \just{$\namer$ separates \texttt{price} from \texttt{fmt}} \\[3pt]
  \texttt{fmt} &: (\ddel, \dupd) \mapsto \text{the update}
    & \just{\cref{def:mergecombine}, \cref{rem:updel}} \\[3pt]
  \NM &= \{\texttt{fmt} \mapsto \texttt{\#f}',\
            \texttt{price} \mapsto \texttt{\#f},\
            \texttt{show} \mapsto \texttt{\#c}'\} &
\end{array}
\]
This merge succeeds, reports nothing, and is wrong in a way no diagnostic covers. The rename has been
undone: \texttt{fmt} is back, because a delete loses to an update. Its other half survived, so
\texttt{price} now names the superseded version of the formatter. Where the ancestor had one
definition the result has two, one of them stale, and \texttt{show} calls the one the user meant to
rename away.

\subsection{An alias does it too, in either of two ways}\label{ex:merge-alias}

An \texttt{alias} also moves names around. What it does depends on whether it is made before or after
the update.

\emph{Before.} $\NA$ aliases the formatter as \texttt{f}, then updates it.
\[
\begin{array}{r@{\;}l@{\quad}l}
  \NA &= \{\texttt{fmt} \mapsto \texttt{\#f}',\ \texttt{f} \mapsto \texttt{\#f},\
            \texttt{show} \mapsto \texttt{\#c}\}
    & \just{\texttt{show} is the ancestor's} \\[3pt]
  \render(\ppe, \texttt{show}) &= ``\texttt{show = }\ulam{n}{\uapp{\uvar{\texttt{f}}}{\uvar n}}\text{''}
    & \just{the dependency still has a name} \\[3pt]
  \uparse(\cdot) &\rightsquigarrow \uref{\texttt{\#f}}
    & \just{and it is the old one}
\end{array}
\]
\texttt{show} is selected as a dependent and rebuilt, and the rebuild returns it unchanged.
\Cref{def:update} rebuilds through text: the renderer shows the dependency under a name, the alias is
a name the old definition still has, and parsing it back arrives where it started. The caller is
pinned. \Cref{prop:follow-only} says this cannot be requested, but it can be arranged.

\emph{After.} $\NA$ updates the formatter, letting \texttt{show} rebuild, and only then aliases it as
\texttt{a}. $\NB$ edits \texttt{show}.
\[
\begin{array}{r@{\;}l@{\quad}l}
  \namer(\texttt{\#f}') &= \texttt{a}
    & \just{fewest segments, then alphabetical} \\[3pt]
  \namer(\texttt{\#f}) &= \texttt{fmt}
    & \just{$\NA$ has no name for it; $\NB$ does} \\[3pt]
  \delta_{\NA}(\texttt{show}) &= \dupd
    & \just{$\ulam{n}{\uapp{\uvar{\texttt{fmt}}}{\uvar n}}$ became $\ulam{n}{\uapp{\uvar{\texttt{a}}}{\uvar n}}$} \\[3pt]
  \texttt{show} &\text{ is conflicted} &
\end{array}
\]
This is \cref{ex:merge-rename} with an alias in place of a \texttt{move}, and it arrives the same way:
$\NA$'s half of the conflict file is the ancestor's own definition, printed under a name $\NA$ never
applied to it directly. No name was taken away: \texttt{fmt} still denotes what $\NA$ means by it,
and every ancestor name is still bound. Adding one was enough. \Cref{cor:noedit-user} excludes
aliases for this reason, and \cref{rem:onlyrenames} states the contrapositive over names, not renames.

The four examples exhaust the ways a change of names meets a concurrent edit, and none behaves as it
should. \Cref{ex:merge-rename} reports a conflict that is not one, \cref{ex:merge-dangling} fails with
no conflict, \cref{ex:merge-dup} succeeds and corrupts the namespace, and \cref{ex:merge-alias} either
does the first or quietly defeats propagation, depending on when the alias was made. The common cause
is \cref{def:mergediff}: a diff of two maps keyed by name cannot see a definition's name change.


% =====================================================================
\section{Comparison with arbor}\label{sec:contrast}

This section compares the design with \emph{arbor}, a content-addressed computational substrate we
are developing. A companion formalization exists in the same style: \arborcore{}, an untyped
$\lambda$-calculus over a store, a namespace and an evaluation cache, with an edit calculus; and
\arborstlc{}, a simply-typed extension. Readers who came for the Unison model can stop here.

arbor is a substrate intended to host several languages, not a language. It shares Unison's core
commitments (definitions identified by hash, names as metadata, immutability, derived data keyed by
hash) and diverges on most of the mechanisms below. Its formalization is machine-checked in Agda,
which is why several rows say ``proved''. That is a statement about arbor's much smaller object
language, not a claim that anything about Unison is provable at the same cost.

\begin{center}
\small
\begin{tabular}{@{}p{0.17\textwidth}p{0.38\textwidth}p{0.38\textwidth}@{}}
\toprule
& \textbf{Unison} & \textbf{arbor} \\
\midrule
Store key
& $\Hash \pfn \Comp$; a definition is $\UDer{h}{i}$ (\cref{def:ustore})
& $\Hash \pfn \Node$: one shallow node per hash, children by hash. A definition is a hash
  (\coreref{def:store}) \\
\addlinespace
Mutual recursion
& In the hasher. Presentation-independence is conditional, and the implementation aborts when the
  condition fails (\cref{thm:ordering}, \cref{rem:notpartial})
& Not yet reached. The leading candidate closes the cycle with a binder in the object language (a
  fixpoint over a product, members as projections), leaving the store's type untouched \\
\addlinespace
Reference transparency
& Yes, at $\Pos = 0$ only (\cref{prop:transparent})
& Opaque in \coreref{def:node}; a decision to adopt transparency was taken in 2026-08 and mechanized
  on a branch, pending merge. Either way there is no positional asymmetry, since there are no
  components \\
\addlinespace
Type
& Inside the term hash; a signature change is an identity change (\cref{def:typeinhash})
& A derived aspect about a hash, observational, not constitutive. Forced by wanting a store that can
  hold ill-typed and partially written terms \\
\addlinespace
Primitives
& $\UBuiltin{b}$, named by text; extrinsic (\cref{rem:builtin})
& Manually versioned procedure tags; also extrinsic, by a different mechanism \\
\addlinespace
Namespace
& A relation; conflict representable, unconflictedness demanded by the operations that need it
  (\cref{rem:conflict})
& A partial function; multiplicity pushed out into a family of namespaces, one per branch or agent \\
\addlinespace
Cache
& Keyed by a reference that need not denote anything stored (\cref{prop:watchdom}); no invalidation
  (\cref{cor:wnoinval})
& The same, including the argument (\coreref{cor:cache}). Neither system attaches an
  \emph{is-a-definition} status to stored data that the cache would have to respect \\
\addlinespace
Update mechanism
& Render $\to$ parse $\to$ typecheck $\to$ re-add (\cref{def:update})
& A structural rewrite over stored nodes. No surface syntax is involved, so no round-trip obligation
  arises \\
\addlinespace
Propagation
& One strategy, scope fixed (\cref{prop:follow-only})
& Pin / follow / explicit over a scope axis. This is richer than the closest comparable system ships,
  so it should not be read as describing prior art \\
\addlinespace
Update failure
& Evict the broken names, hand back source (\cref{rem:residual})
& Commit and report a residual set of names \\
\addlinespace
Round-trip lemma
& Load-bearing, unstated, maintained by a golden corpus (\cref{rem:printparse})
& \coreref{prop:print-elab}, proved, and not depended upon \\
\addlinespace
Merging two edited namespaces
& Nothing in the store is reconciled (\cref{prop:merge-store}); all of it happens in $\UNames$
& The same, for the same reason. An immutable store buys less here than the usual summary suggests \\
\addlinespace
Telling an authored edit from a propagated one
& The syntactic hash: did it \emph{render} differently (\cref{def:synhash})
& No design recorded. arbor has the reverse index a different answer would need, and has not been
  asked the question \\
\bottomrule
\end{tabular}
\end{center}

\subsection*{Three observations we take from the comparison}

\textbf{On the round-trip property.} \Cref{prop:print-parse} identifies a correctness obligation that
a shipping system carries and maintains by testing. arbor has the corresponding theorem proved, for
its own much smaller languages, and does not depend on it, since its migration is a structural
rewrite that never goes through surface syntax. The observation concerns the shape of the guarantee,
not either system's rigor: one system depends on a property it does not state, and the other states a
property it does not depend on. A specification of the pretty-print environment under which the
property becomes provable at Unison's scale would be useful, and is out of reach here.

\textbf{On the failure case.} \Cref{rem:residual} is the row where we think Unison is ahead. arbor's
migration commits and reports a residual, a set of names that no longer typecheck. Unison's answer is
more useful: the broken definitions leave the namespace and come back as code in a file, ready to
edit. The namespace never holds a broken state, and the report is the work itself.

\textbf{On where the difficulty lives.} \Cref{prop:merge-store}, the property a content-addressed
codebase is usually sold on, is true but is nearly all that immutability buys at merge time. Every
hard case in \cref{sec:umerge} is in $\UNames$, a relation over names people chose with no content
addressing in it. The mechanism Unison built to cope, the syntactic hash, puts the names back into a
hash (\cref{def:synhash}). The design removes names from identity; its merge algorithm restores them
for the duration of one comparison. We do not read that as incoherent, and \cref{lem:prop-silent}
shows why it is the right instinct. It does locate the residual difficulty, and a system inheriting
the store model should not expect an easier merge to come with it.

% =====================================================================
\appendix

% =====================================================================
\section{Provenance, and how to check this}\label{sec:paths}

Everything here is modeled on Unison at commit \texttt{db60ce2} (2026-08-07). Re-check the pin with
\texttt{git -C <unison> rev-parse HEAD}. If it has moved, the citations below are still true at that
commit; they are dated, not necessarily wrong.

The arithmetic in this document has also been executed. A reimplementation of
\crefrange{def:tokens}{def:update} over this object language exists, and each statement in
\crefrange{sec:ulang}{sec:uupdate} is a test in it that names the statement it checks. \Cref{sec:umerge}
is not yet part of that reimplementation; its evidence is described below. The reimplementation is
not Unison and cannot settle a question about Unison's intent. It settles whether our definitions say
what we claim. It changed three statements: \cref{rem:transparency-limit} (confirmed, and at $k = 1$),
\cref{thm:ordering}(b) (which had to split, because renaming a tied member does not always move
$\comphash$), and \cref{prop:print-parse} (whose first implementation failed the property on a printer
that dropped parentheses around an application in argument position, the class of defect the golden
corpus of \cref{rem:printparse} exists to catch).

Some findings rest on direct observation of \texttt{ucm} 1.1.1 instead of on the model.
\Cref{rem:innercycle}'s three hashes are what \texttt{ucm} reports for the three definitions as
written, and the top-level abort is the tool refusing the same pair of bindings lifted out of the
block. The model reproduces the top-level behaviour but not yet the inner one. Each of
\crefrange{ex:merge-rename}{ex:merge-alias} was run as a transcript over two or three definitions,
and the outcomes reported (a conflict, a no-conflict failure naming a deleted identifier, a clean
merge with the rename undone, and an alias that either pins a dependent or conflicts it) are what the
tool produced, not what the model predicts. \Cref{prop:merge-asym} was obtained by running one pair
of branches in both directions, and \cref{prop:follow-only}'s pinning case by running one
\texttt{update} with and without a preceding \texttt{alias} and comparing the dependent's hash. This
evidence does not show that our definitions are internally coherent. It does show that the behavior
is Unison's and not an artifact of our reading.

Citations are given where a claim is surprising, load-bearing, or likely to be doubted, and are not
exhaustive. This is an idealized model, not a transcription, and a tag on every definition would
suggest a closer correspondence than the one claimed. A margin tag has the form \texttt{key:lines},
with keys as follows.

\begin{center}
\small
\begin{tabular}{@{}ll@{}}
\toprule
Key & File \\
\midrule
\texttt{abt}      & \texttt{unison-hashing-v2/src/Unison/Hashing/V2/ABT.hs} \\
\texttt{term}     & \texttt{unison-hashing-v2/src/Unison/Hashing/V2/Term.hs} \\
\texttt{type}     & \texttt{unison-hashing-v2/src/Unison/Hashing/V2/Type.hs} \\
\texttt{branchh}  & \texttt{unison-hashing-v2/src/Unison/Hashing/V2/Branch.hs} \\
\texttt{tokenizable} & \texttt{unison-hashing-v2/src/Unison/Hashing/V2/Tokenizable.hs} \\
\texttt{ref}      & \texttt{codebase2/core/U/Codebase/Reference.hs} \\
\texttt{referent} & \texttt{codebase2/core/U/Codebase/Referent.hs} \\
\texttt{cterm}    & \texttt{codebase2/codebase/U/Codebase/Term.hs} \\
\texttt{ctype}    & \texttt{parser-typechecker/src/Unison/Codebase/Type.hs} \\
\texttt{tformat}  & \texttt{codebase2/codebase-sqlite/U/Codebase/Sqlite/Term/Format.hs} \\
\texttt{queries}  & \texttt{codebase2/codebase-sqlite/U/Codebase/Sqlite/Queries.hs} \\
\texttt{branchv3} & \texttt{codebase2/codebase/U/Codebase/BranchV3.hs} \\
\texttt{ops}      & \texttt{codebase2/codebase-sqlite/U/Codebase/Sqlite/Operations.hs} \\
\texttt{sql}      & \texttt{codebase2/codebase-sqlite/sql/create.sql} \\
\texttt{ufile}    & \texttt{parser-typechecker/src/Unison/UnisonFile.hs} \\
\texttt{update}   & \texttt{unison-cli/src/Unison/Codebase/Editor/HandleInput/Update2.hs} \\
\texttt{load}     & \texttt{unison-cli/src/Unison/Codebase/Editor/HandleInput/Load.hs} \\
\texttt{main}     & \texttt{unison-cli/src/Unison/Main.hs} \\
\texttt{decl}     & \texttt{codebase2/codebase/U/Codebase/Decl.hs} \\
\texttt{termhash} & \texttt{codebase2/codebase-sqlite-hashing-v2/src/U/Codebase/Term/Hashing.hs} \\
\texttt{tparse}   & \texttt{parser-typechecker/src/Unison/Syntax/TermParser.hs} \\
\texttt{fparse}   & \texttt{parser-typechecker/src/Unison/Syntax/FileParser.hs} \\
\texttt{coreterm} & \texttt{unison-core/src/Unison/Term.hs} \\
\texttt{branchtype} & \texttt{parser-typechecker/src/Unison/Codebase/Branch/Type.hs} \\
\texttt{convert}  & \texttt{parser-typechecker/src/Unison/Hashing/V2/Convert.hs} \\
\texttt{inputpat} & \texttt{unison-cli/src/Unison/CommandLine/InputPatterns.hs} \\
\texttt{names}    & \texttt{unison-core/src/Unison/Names.hs} \\
\texttt{ppenames} & \texttt{parser-typechecker/src/Unison/PrettyPrintEnv/Names.hs} \\
\texttt{msynhash} & \texttt{unison-merge/src/Unison/Merge/Synhash.hs} \\
\texttt{msynhashed} & \texttt{unison-merge/src/Unison/Merge/Synhashed.hs} \\
\texttt{mdiff}    & \texttt{unison-merge/src/Unison/Merge/Diff.hs} \\
\texttt{mdiffblob} & \texttt{unison-merge/src/Unison/Merge/Diffblob.hs} \\
\texttt{mtwodiff} & \texttt{unison-merge/src/Unison/Merge/TwoDiffOps.hs} \\
\texttt{mcombine} & \texttt{unison-merge/src/Unison/Merge/CombineDiffs.hs} \\
\texttt{mnarrow}  & \texttt{unison-merge/src/Unison/Merge/Narrow.hs} \\
\texttt{mrename}  & \texttt{unison-merge/src/Unison/Merge/Rename.hs} \\
\texttt{mrender}  & \texttt{unison-merge/src/Unison/Merge/Render.hs} \\
\texttt{mblob}    & \texttt{unison-merge/src/Unison/Merge/Mergeblob.hs} \\
\texttt{merge2}   & \texttt{unison-cli/src/Unison/Codebase/Editor/HandleInput/Merge2.hs} \\
\bottomrule
\end{tabular}
\end{center}

Unison has no publication, so the citable artifact for any claim about its behavior is the source at
a pinned commit. We encourage anyone building on this document to cite it that way.

Corrections are welcome, particularly from Unison's authors. The code underdetermines the intent at
two points, where we are correspondingly least confident: whether the position-$0$ boundary of
\cref{prop:transparent} is intended, and whether the two orderings of \cref{rem:twoorders} are meant
to be distinct. Both are questions about intent, not about what the code does.

The question we would most like answered is whether the interaction between renaming and
\cref{def:mergediff} is regarded as a defect. The rename information is computed and unused
(\cref{cor:unchecked}), and the one case where the misclassification was serious enough to patch
around is \cref{rem:libcanon}, so the phenomenon is known. We cannot tell from the source whether
\cref{ex:merge-dup}, a clean merge that undoes a rename and leaves a stale alias, is understood to
follow from the same cause.

One claim contradicts a frequently repeated description of the system: \cref{rem:transparency-limit},
that a definition's \emph{component} hash is not its \emph{term} hash. The familiar illustration of
reference transparency therefore does not hold for whole definitions, though the term-level property
does. The construction that satisfies the stronger reading, \texttt{hashClosedTerm}
\src{term:124-125}, is used for anonymous and evaluated expressions and not for storing a named
definition, which makes the difference observable. We would welcome confirmation, but are no longer
unsure: an independent implementation of \cref{def:comphash} reproduces the gap at $k = 1$, where it
is easiest to check by hand.

\end{document}
